\begin{textsample}{POS Dim 3 – 2018\_c6 – Score 3.00 – 2018\_c6/enade\_2018\_ci\_cont\_3fded95eb1.txt}  \label{ex:f3_pos_111}
CAPÍTULO 6 ANÁLISE TÉCNICA DA PROVA Este capítulo tem por objetivo apresentar o desempenho dos estudantes concluintes de Ciências Contábeis no Enade/2018.
Para isso, foram calculadas as estatísticas básicas da prova em seu todo ( seção 6.1.1 ), bem como as estatísticas relacionadas aos Componentes de Formação Geral ( seção 6.1.2 ) e de Conhecimento Específico da Área ( seção 6.1.3 ).
Dadas as suas características, foram analisadas, em separado, as questões objetivas ( seção 6.2 ) e as questões discursivas ( seção 6.3 ).
Para as questões objetivas, foram disponibilizados os índices de facilidade e de discriminação Ponto-Bisserial, também em separado, para os Componentes de Formação Geral ( seção 6.2.1 ) e de Conhecimento Específico ( seção 6.2.2 ).
De cada componente, uma das questões foi escolhida para exemplificar a análise gráfica, relacionando as alternativas escolhidas pelos estudantes ( inclusive o gabarito ) com o número de acertos no componente.
No Anexo I, apresenta -se a íntegra da análise gráfica para todas as questões objetivas.
Para cada uma das questões discursivas, os conteúdos dos tipos mais comuns de respostas dos estudantes são apresentados e comparados com o padrão de respostas esperado ( ver Anexo VIII com o padrão de respostas ).
Tomando -se como base as duas questões discursivas do Componente de Formação Geral, na seção 6.3.1.6, são apresentados comentários sobre a correção das respostas em relação à Língua Portuguesa.
Nas tabelas, constam as seguintes estatísticas das notas20 : média do desempenho na prova, erro padrão da média, desvio padrão, nota mínima, mediana e nota máxima para cada um de seus componentes.
Tais estatísticas contemplam o total de estudantes concluintes da Área de Ciências Contábeis inscritos e presentes à prova do Enade/2018, tendo em vista agregações, ou por Grandes Regiões ou o país como um todo, ou por Categoria Administrativa, Organização Acadêmica da IES e Modalidade de Ensino.
Em relação aos gráficos de distribuição de notas, o intervalo considerado foi de 10 unidades, aberto à esquerda e fechado à direita, com exceção do primeiro intervalo, [ 0 ; 10 ], fechado em ambos os extremos.
Para os gráficos de distribuição das notas das questões discursivas, foram consideradas mais duas categorias : questão em branco21 e nota zero. 20 Uma definição dessas estatísticas pode ser encontrada no Glossário. 21 Nesse grupo estão incluídas também as respostas classificadas como nulas ou desconsideradas. 6.1 ESTATÍSTICAS BÁSICAS DA PROVA Esta seção apresenta estatísticas selecionadas e histogramas da nota geral ( 6.1.1 ) e de cada componente : Formação Geral ( 6.1.2 ) e Conhecimento Específico ( 6.1.3 ).
São, também apresentadas estatísticas selecionadas de subpopulações, caracterizadas por Grande Região, Categoria Administrativa, Organização Acadêmica e Modalidade de Ensino. 6.1.1 Estatísticas Básicas Gerais Na Tabela 6.1, são apresentadas as Estatísticas Básicas da prova, por Grande Região, dos estudantes concluintes de Ciências Contábeis.
A Média das notas da prova, como um todo ( nas seções seguintes serão analisados os Componentes de Formação Geral e de Conhecimento Específico ), foi 36,7, sendo que os alunos da região Norte obtiveram a Média mais baixa ( 34,6 ), e os da região Sudeste obtiveram a Média mais alta ( 38,2 ).
As demais médias foram : 36,1, na região Nordeste, 35,8, na região Sul e 35,7, na região Centro-Oeste.
O Desvio padrão para o Brasil, como um todo, foi 12,5, sendo o maior Desvio padrão encontrado na região Sudeste ( 12,7 ), e o menor, na região Norte ( 11,1 ), indicando uma dispersão um pouco menor das notas dessa última região.
A região que obteve a maior nota Máxima foi a Centro-Oeste ( 90,1 ), ao passo que a região que atingiu a menor nota Máxima foi a Norte ( 76,5 ).
A Mediana do Brasil, como um todo, foi 35,7, sendo a maior Mediana obtida na região Sudeste ( 37,2 ), e a menor obtida na Norte ( 33,7 ).
A nota Mínima do Brasil foi obtida em todas as regiões, onde esta foi zero.
Considerando -se as notas segundo Grande Região, observa -se que não existe diferença estatisticamente significativa, ao nível de 95 \%, entre as regiões Nordeste, Sul e Centro-Oeste.
Entre as regiões Norte ( menor Média ) e a Sudeste ( maior Média ), existe diferença estatisticamente significativa, bem como destas com as demais regiões²². ²² Todas as comparações deste capítulo utilizam os intervalos de 95 \%.
Os erros-padrão da média que possibilitam os testes estão disponíveis nas tabelas.
O comportamento das notas dos estudantes de todo o Brasil pode ser observado no Gráfico 6.1, que apresenta um histograma com a distribuição das mesmas : uma distribuição unimodal com a moda no intervalo ( 30 ; 40 ], seguido pelo intervalo ( 20 ; 30 ].
Na Tabela 6.2, são apresentadas informações referentes à Média da nota final, desagregadas por Categoria Administrativa, por Organização Acadêmica e por Modalidade de Ensino.
A Média das notas da prova, como um todo, foi 36,7.
Em relação à Categoria Administrativa, os estudantes das IES Públicas obtiveram Média mais alta ( 40,4 ) que a Média nacional, e os das IES Privadas obtiveram Média mais baixa ( 36,0 ).
Observa -se que existe diferença estatisticamente significativa entre as médias das notas das IES Públicas e as das Privadas.
A diferença entre as médias das regiões Sudeste e Norte ( 3,6 ), a maior e a menor Média, é inferior à diferença entre a das IES Públicas e a das Privadas ( 4,4 ), caracterizando -se uma maior diversidade administrativa do que regional.
O Desvio padrão para as IES Públicas ( 13,6 ) foi superior ao do Brasil, como um todo ( 12,5 ), indicando uma dispersão um pouco maior das notas nesta Categoria Administrativa.
No tocante à Organização Acadêmica, os CEFET/IFET, as Universidades e os Centros Universitários obtiveram Média mais alta que a nacional ( 39,8, 37,4 e 36,8, respectivamente ).
A Média das Faculdades foi menor do que a nacional ( 35,3 ).
Constata -se que existe diferença estatisticamente significativa, ao nível de 95 \%, entre a Média mais baixa das Faculdades e a das demais Organizações Acadêmicas.
A Média da Modalidade Educação a Distância ( 34,2 ) foi inferior à Média da Educação Presencial ( 37,5 ), e a diferença entre elas é estatisticamente significativa, ao nível de 95 \%. 6.1.2 Estatísticas Básicas no Componente de Formação Geral Na Tabela 6.3, são apresentadas as Estatísticas Básicas em relação ao componente da prova que avalia a Formação Geral dos estudantes concluintes.
Os alunos de todo o Brasil obtiveram desempenho médio de 41,3.
Quanto à variabilidade, o Desvio padrão das notas dos estudantes do Brasil, como um todo, foi 16,5.
A maior Média foi obtida na região Sudeste ( 42,0 ), e a menor, na região Centro-Oeste ( 40,1 ).
As demais médias foram : 40,2, na região Norte, 41,4, na região Nordeste e 41,0, na região Sul.
Já o maior Desvio padrão foi obtido nas regiões Nordeste e Centro-Oeste ( 16,8 ), e o menor, nas regiões Norte e Sul ( 16,1 ) ; o Desvio padrão da região Sudeste foi 16,6.
A maior nota no Componente de Formação Geral da prova do Enade/2018 foi 94,4, obtida por, pelo menos, um aluno na região Sudeste.
A menor nota Máxima foi obtida na região Norte ( 89,6 ).
A Mediana do Brasil, como um todo, foi 40,8, sendo a menor Mediana encontrada na região Centro-Oeste ( 39,3 ), e a maior, na região Sudeste ( 41,6 ).
A nota Mínima nessa parte foi zero em todas as regiões.
Considerando -se as notas, segundo Grande Região, observa -se que existe diferença estatisticamente significativa entre a Média das notas do Componente de Formação Geral da região Sudeste ( a mais alta ) e as demais regiões.
No Gráfico 6.2, é possível verificar a avaliação do desempenho dos estudantes do Componente de Formação Geral, a partir do histograma da distribuição das notas correspondentes.
A distribuição é unimodal, com moda em ( 40 ; 50 ], intervalo imediatamente seguinte ao modal da distribuição de notas da prova, como um todo ( Gráfico 6.1 ).
Nota -se, ainda, que, no Gráfico 6.2, as notas apresentam maior dispersão do que as no Gráfico 6.1 ( distribuição das notas da prova ), confirmado pela comparação dos desvios padrão : 12,5 para a nota da prova, como um todo, e 16,5 para o Componente de Formação Geral.
Vale a pena ressaltar que o intervalo anterior ao modal, ( 30 ; 40 ], tem frequência muito próxima ao modal.
Na Tabela 6.4, são apresentadas as informações referentes ao desempenho dos concluintes do Componente de Formação Geral, em diferentes agregações : Categoria Administrativa, Organização Acadêmica e Modalidade de Ensino.
Observa -se que existe diferença estatisticamente significativa entre as médias dos tipos de Categoria Administrativa : a maior Média obtida por estudantes de IES Públicas ( 46,5 ) e a menor, pelos de IES Privadas ( 40,4 ).
Nota -se que, também, há diferença estatisticamente significativa entre a menor das médias dos quatro tipos de Organização Acadêmica e as demais, a menor para as Faculdades ( 39,3 ), para Universidades ( 42,7 ), Centros Universitários ( 39,3 ) e a maior para os CEFET/IFET ( 44,9 ).
No que se refere à Modalidade de Ensino, a Educação Presencial apresentou maior Média do que a Educação a Distância ( médias 41,6 e 40,6, respectivamente ).
Tais médias apresentam, também, diferença estatisticamente significativa. 6.1.3 - Estatísticas Básicas do Componente de Conhecimento Específico Na Tabela 6.5, são apresentadas as Estatísticas Básicas referentes ao Componente de Conhecimento Específico da área de Ciências Contábeis.
A Média do desempenho dos alunos do Brasil, como um todo, foi 35,1.
A maior Média foi obtida na região Sudeste ( 36,9 ), e a menor, na região Norte ( 32,7 ).
As demais médias foram : 34,3, nas regiões Nordeste e Centro-Oeste e 34,1, na região Sul.
Quanto à variabilidade das notas, o Desvio padrão do Brasil, como um todo, foi 13,8, sendo o maior Desvio padrão observado na região Sudeste ( 14,1 ) e o menor, na região Norte ( 12,4 ).
Os demais desvios foram : 13,7, nas regiões Nordeste e Centro-Oeste e 13,4, na região Sul.
A Mediana das notas dos estudantes de todo o Brasil foi 33,9.
A maior Mediana ocorreu na região Sudeste ( 36,9 ), e a menor, na região Norte ( 31,9 ).
As demais medianas foram : 32,9, na região Nordeste, 32,6, na região Sul e 32,4, na região Centro-Oeste.
A nota Máxima do Brasil, como um todo, foi 94,2, sendo obtida por, pelo menos, um aluno na região Centro-Oeste.
As demais notas máximas foram : 77,5, nas regiões Norte ( a menor ), 87,1, na região Nordeste, 92,7, na região Sudeste e 89,5, na região Sul.
A nota Mínima foi zero em todas as regiões.
Observa -se que não existe diferença estatisticamente significativa entre as Médias, das notas do Componente de Conhecimento Específico, das regiões Nordeste, Sul e Centro-Oeste.
Entre as regiões Norte ( menor Média ) e a Sudeste ( maior Média ) existe diferença estatisticamente significativa, bem como destas e as demais regiões23. 23 Todas as comparações deste capítulo utilizam os intervalos de 95 \%.
Os erros-padrão da média que possibilitam os testes estão disponíveis nas tabelas.
Assim como nos Gráficos 6.1 e 6.2, no Gráfico 6.3, apresentado a seguir, é possível ser feita uma avaliação do desempenho de concluintes de Ciências Contábeis, em relação ao Componente de Conhecimento Específico, com um histograma da distribuição das notas correspondentes.
Esta também é uma distribuição unimodal, e o grupo modal é o ( 30 ; 40 ], mesmo grupo modal da prova, como um todo, e abaixo do grupo modal para a Formação Geral.
Na Tabela 6.6, é apresentada uma comparação dos resultados em relação à Categoria Administrativa, à Organização Acadêmica e à Modalidade de Ensino, agora levando -se em conta o desempenho de alunos do Componente de Conhecimento Específico da prova.
No que se refere à Organização Acadêmica, a maior Média foi a dos CEFET/IFET ( 38,1 ), vindo a seguir a das Universidades ( 35,6 ), a dos Centros Universitários ( 35,5 ) e, por fim, a das Faculdades ( 33,9 ).
O maior Desvio padrão, e acima do valor para o Brasil, como um todo ( 13,8 ), foi o das Universidades ( 14,2 ).
As Universidades obtiveram a maior nota Máxima ( 94,2 ).
Os Centros Universitários obtiveram nota Máxima 92,7 ; as Faculdades obtiveram nota Máxima 90,7 ; e os CEFET/IFET, 71,7, a menor.
As medianas foram : 34,4, nas Universidades e nos Centros Universitários ; 32,6, nas Faculdades, a menor delas ; e 37,2, nos CEFET/IFET, a maior.
A nota Mínima foi zero para quase todas as Organizações Acadêmicas, com exceção dos CEFET/IFET ( 10,9 ).
Observa -se que existe diferença estatisticamente significativa, ao nível de 95 \%, no Componente de Conhecimento Específico entre as notas das Faculdades ( a menor Média ) e as das demais Organizações Acadêmicas.
Quanto à Categoria Administrativa, observa -se um comportamento semelhante àquele da parte de Formação Geral e à prova, como um todo, ou seja, existe diferença estatisticamente significativa entre as médias das IES Públicas ( 38,4 ) e as das IES Privadas ( 34,5 ).
Neste caso, também, a maior Média foi obtida por alunos de IES Públicas de ensino.
Quanto ao comportamento das médias, no que toca à Modalidade de Ensino, há diferença estatisticamente significativa entre a Média da Educação Presencial ( 36,1 ) e a da Educação a Distância ( 32,0 ), sendo a maior Média, também, para os cursos presenciais. 6.2 ANÁLISE DAS QUESTÕES OBJETIVAS Esta seção apresenta estatísticas selecionadas e histogramas das Questões Objetivas de Formação Geral ( 6.2.1 ) e de Conhecimento Específico ( 6.2.2 ).
São, também, apresentadas e comparadas as médias das subpopulações caracterizadas por Grande Região. 6.2.1 Componente de Formação Geral Na Tabela 6.7, são apresentadas as Estatísticas Básicas relativas às oito questões objetivas do componente da prova que abrange a Formação Geral dos estudantes.
A Média do Brasil foi 45,9.
A menor Média foi encontrada na região Norte ( 43,6 ), e a maior, na região Sudeste ( 46,9 ).
As demais médias foram : 46,3, na região Nordeste, 44,7, na região Sul e 45,4, na região Centro-Oeste.
O Desvio padrão do Brasil foi 22,1, sendo o maior Desvio padrão encontrado na região Centro-Oeste ( 22,4 ), e o menor, na região Norte ( 21,3 ).
Os demais desvios foram : 22,3, na região Nordeste, 22,2, na região Sudeste e 21,8, na região Sul.
As medianas do Brasil, como um todo, e de duas das cinco regiões foram 50,0 - as regiões Norte, Sul e Centro-Oeste com Mediana 37,5.
A nota Máxima 100,0 foi alcançada em todas as Grandes Regiões.
As notas mínimas ( 0,0 ) foram iguais para todas as regiões.
Na Tabela 6.8, são apresentados o Índice de Facilidade e o Índice de Discriminação ( Ponto-Bisserial ) para cada uma das questões objetivas do Componente de Formação Geral.
Quanto ao Índice de Facilidade, foram usadas as seguintes cores para diferenciar o nível de dificuldade da questão : • Azul para as questões classificadas com índice Muito fácil ( \&gt ; =0,86 ), verde para as questões classificadas com índice Fácil ( 0,61 a 0,85 ), amarelo para as questões classificadas com Médio ( 0,41 a 0,60 ), vermelho para as questões classificadas com Difícil ( 0,16 a 0,40 ) e roxo para as questões classificadas com Muito difícil ( \&lt ; =0,15 ).
Na Tabela 6.9, é apresentada a distribuição das questões, levando -se em conta, simultaneamente, a classificação dos dois índices.
As oito questões obtiveram Índice de Discriminação Muito bom e figuraram entre três níveis de dificuldade Fácil, Médio e Difícil : uma classificada na categoria Fácil ( a questão 2 ) do Índice de Facilidade, cinco, na categoria Médio ( questões 1, 3, 6, 7 e 8 ), e duas, na categoria Difícil ( questões 4 e 5 ).
Em particular, a questão 1 foi a que apresentou o maior poder discriminatório, com índice 0,49, e foi considerada Médio em termos de facilidade, com uma proporção de 49 \% de acertos.
O máximo de acertos foi alcançado pela questão 2 com um Índice de Facilidade de 0,62.
A questão 5 foi considerada Difícil, com Índice de Facilidade 0,37 e, por sua vez, o Índice de Discriminação foi Muito bom ( 0,46 ).
A questão de número 4 apresentou um Índice de Já quanto ao Índice de Discriminação, foram usadas as seguintes cores para qualificar a questão : • As questões classificadas com índice Fraco receberam a cor vermelha ( \&lt ; =0,19 ), as classificadas com Médio receberam a cor amarela ( 0,20 a 0,29 ), as classificadas com Bom receberam a cor verde ( 0,30 a 0,39 ) e as classificadas com Muito bom ( \&gt ; =0,40 ) receberam a cor azul.
As questões objetivas do Componente de Formação Geral, segundo o Índice de Facilidade, foram assim avaliadas : das oito questões, nenhuma teve o Índice de Facilidade classificado como Muito fácil, e uma questão foi tida como Fácil, por ter índice de acertos 0,62.
Cinco questões foram consideradas com Índice de Facilidade Médio, situando -se no intervalo entre 0,41 e 0,60, ou seja, houve entre 41,0 \% e 60,0 \% de acertos, enquanto duas questões foram classificadas na categoria Difícil, situando -se no intervalo entre 0,16 e 0,40.
Nenhuma questão apresentou menos de 15 \% de acertos, razão pela qual não houve questão classificada como Muito difícil.
O Índice de Facilidade variou de 0,24 a 0,62, e o de Discriminação, de 0,41 a 0,49.
Na Tabela 6.9, é apresentada a distribuição das questões, levando -se em conta, simultaneamente, a classificação dos dois índices.
As oito questões obtiveram Índice de Discriminação Muito bom e figuraram entre três níveis de dificuldade Fácil, Médio e Difícil : uma classificada na categoria Fácil ( a questão 2 ) do Índice de Facilidade, cinco, na categoria Médio ( questões 1, 3, 6, 7 e 8 ), e duas, na categoria Difícil ( questões 4 e 5 ).
Em particular, a questão 1 foi a que apresentou o maior poder discriminatório, com índice 0,49, e foi considerada Médio em termos de facilidade, com uma proporção de 49 \% de acertos.
O máximo de acertos foi alcançado pela questão 2 com um Índice de Facilidade de 0,62.
A questão 5 foi considerada Difícil, com Índice de Facilidade 0,37 e, por sua vez, o Índice de Discriminação foi Muito bom ( 0,46 ).
A questão de número 4 apresentou um Índice de Facilidade de 0,24 ( Difícil ), ou seja, menos do que um quarto dos estudantes conseguiu resolvê -la, dentro do universo de participantes.
Seu Índice de Discriminação foi Muito bom ( 0,46 ).
No Gráfico 6.4, para exemplificar, analisa -se o comportamento da questão de número 4 de Formação Geral.
Trata -se de uma questão considerada Difícil, em relação à facilidade e a que obteve o quarto maior Índice de Discriminação dessa parte da prova.
Neste gráfico, cada uma das cinco curvas representa o percentual de respostas em determinada alternativa da questão, em função do número de acertos dos estudantes nessa parte da prova ( Formação Geral/Múltipla Escolha ), antes de possíveis eliminações pelo critério do Ponto-Bisserial.
A curva em verde corresponde à alternativa B, a alternativa correta para essa questão.
Observa -se que entre os estudantes com menor número de acertos, nessa parte do exame, a situação mais frequente foi a escolha da alternativa C ( em preto ), incorreta.
Por exemplo, entre os estudantes que acertaram três questões, 40,9 \% escolheram a alternativa C ( em preto ), 18,9 \% escolheram a alternativa E ( em vermelho ), 16,2 \% escolheram a alternativa D ( em \textbf{laranja} ), 13,5 \% escolheram a alternativa B ( em verde ) e 10,5 \% escolheram a alternativa A ( azul ).
Dentre os que acertaram três respostas entre as questões de múltipla escolha de Formação Geral, 0,0 \% deixou a questão em branco e 0,1 \% marcou mais de uma alternativa, invalidando a questão.
À medida que o número de acertos aumenta, indicando desempenho melhor nessa parte da prova, aumenta, concomitantemente, a proporção de estudantes que selecionaram a alternativa correta B, atingindo 100 \% para os estudantes com oito acertos.
Essa análise permite verificar como a questão discriminou os grupos de desempenho, justificando -se o alto índice obtido na questão.
Cumpre notar que não é possível inferir deste gráfico nem o índice de facilidade ( que seria uma média da proporção ponderada pela quantidade de alunos com cada uma das notas ), nem o índice de discriminação Ponto-Bisserial, por razão equivalente.
No caso extremo, no qual a grande concentração dos acertos dos alunos fosse abaixo de quatro, o Facilidade de 0,24 ( Difícil ), ou seja, menos do que um quarto dos estudantes conseguiu resolvê -la, dentro do universo de participantes.
Seu Índice de Discriminação foi Muito bom ( 0,46 ).
No Gráfico 6.4, para exemplificar, analisa -se o comportamento da questão de número 4 de Formação Geral.
Trata -se de uma questão considerada Difícil, em relação à facilidade e a que obteve o quarto maior Índice de Discriminação dessa parte da prova.
Neste gráfico, cada uma das cinco curvas representa o percentual de respostas em determinada alternativa da questão, em função do número de acertos dos estudantes nessa parte da prova ( Formação Geral/Múltipla Escolha ), antes de possíveis eliminações pelo critério do Ponto-Bisserial.
A curva em verde corresponde à alternativa B, a alternativa correta para essa questão.
Observa -se que entre os estudantes com menor número de acertos, nessa parte do exame, a situação mais frequente foi a escolha da alternativa C ( em preto ), incorreta.
Por exemplo, entre os estudantes que acertaram três questões, 40,9 \% escolheram a alternativa C ( em preto ), 18,9 \% escolheram a alternativa E ( em vermelho ), 16,2 \% escolheram a alternativa D ( em \textbf{laranja} ), 13,5 \% escolheram a alternativa B ( em verde ) e 10,5 \% escolheram a alternativa A ( azul ).
Dentre os que acertaram três respostas entre as questões de múltipla escolha de Formação Geral, 0,0 \% deixou a questão em branco e 0,1 \% marcou mais de uma alternativa, invalidando a questão.
À medida que o número de acertos aumenta, indicando desempenho melhor nessa parte da prova, aumenta, concomitantemente, a proporção de estudantes que selecionaram a alternativa correta B, atingindo 100 \% para os estudantes com oito acertos.
Essa análise permite verificar como a questão discriminou os grupos de desempenho, justificando -se o alto índice obtido na questão.
Cumpre notar que não é possível inferir deste gráfico nem o índice de facilidade ( que seria uma média da proporção ponderada pela quantidade de alunos com cada uma das notas ), nem o índice de discriminação Ponto-Bisserial, por razão equivalente.
No caso extremo, no qual a grande concentração dos acertos dos alunos fosse abaixo de quatro, o índice de facilidade seria, obrigatoriamente, abaixo de 20 \% ( neste exemplo ).
Caso a concentração fosse em seis acertos ou mais, o índice seria, obrigatoriamente, acima de 50 \%.
Os gráficos relativos às demais questões de Formação Geral constam do Anexo I. 6.2.2 Componente de Conhecimento Específico Na Tabela 6.10, são apresentadas as Estatísticas Básicas em relação às questões objetivas do Componente de Conhecimento Específico da prova de Ciências Contábeis por Grande Região.
Nessa parte da prova, duas questões foram anuladas pela Comissão Assessora da Área, questões 25 e 26.
Além disso, como será discutido ainda nesta seção, nove questões objetivas do Componente de Conhecimentos Específicos não foram usadas no cômputo das notas por terem sido descartadas pelo critério Ponto-Bisserial.
Assim, as notas foram calculadas com base em 16 das 27 questões objetivas de Conhecimentos Específicos.
A Média do Brasil deste componente foi 38,7.
A menor Média foi observada na região Norte ( 36,6 ), e a maior, na região Sudeste ( 40,4 ).
O Desvio padrão de todo o Brasil foi 14,7, sendo o menor Desvio padrão encontrado na região Norte ( 13,5 ), e o maior, na região Sudeste ( 14,9 ).
A Mediana de todo o Brasil foi 37,5, a mesma encontrada em todas as regiões.
A nota Máxima da prova foi 93,8, obtida nas questões objetivas do Componente de Conhecimento Específico, por, pelo menos, um aluno em quase todas as regiões, a exceção foi a região Norte onde a Máxima foi 87,5.
A nota Mínima foi zero em todas as regiões.
Na Tabela 6.11, são apresentados os Índices de Facilidade e Discriminação ( Ponto-Bisserial ) das questões objetivas do Componente de Conhecimento Específico, para os estudantes de Ciências Contábeis.
Para facilitar a diferenciação das questões, serão usadas as mesmas cores da Tabela 6.8 para as diferentes classificações dos Índices de Facilidade e de Discriminação.
Observa -se que duas questões foram anuladas pela CAA e que as análises dizem respeito às 25 questões válidas.
A partir do Índice de Facilidade obtido, pode -se concluir que três quartos das questões objetivas da prova foram considerados, pelo menos, Difícil : das 25 questões válidas, 18 foram classificadas como Difícil ou como Muito difícil.
Duas questões foram classificadas como Fácil, e outras cinco consideradas como Médio.
Já quanto aos índices de discriminação das questões objetivas do Componente de Conhecimento Específico da prova, obtém -se como resultado a seguinte classificação : sete das 25 questões foram consideradas boas e outras três questões foram classificadas como Muito bom.
Assim, para dois quintos das questões – dez em 25 – os Índices de Discriminação foram Bom ou Muito bom.
Entre as demais, seis delas foram classificadas como Médio, e outras nove, como Fraco, sendo 15, por conseguinte, a quantidade de questões nos dois patamares mais baixos de discriminação.
Constata -se, assim, que a prova – no que se refere ao Componente de Conhecimento Específico – possuía média capacidade de discriminar entre aqueles que dominam ou não o conteúdo.
O Índice de Facilidade variou de 0,09 a 0,79, e o de Discriminação, de 0,01 a 0,42.
Na Tabela 6.12, é apresentada a distribuição das questões levando -se em conta, simultaneamente, a classificação dos dois índices.
Entre as questões que alcançaram os maiores índices de discriminação, com Índice de Discriminação Bom e Muito bom, a questão de número 32 teve índice 0,42, e, quanto ao índice de Facilidade, foi classificada como Médio, com 44 \% dos estudantes marcando a opção correta.
A questão de número 10 foi a mais difícil entre as 25 questões específicas válidas, com baixo Índice de Facilidade de apenas 9 \% de acertos.
Essa questão apresentou poder discriminatório muito baixo, 0,05, o que comprova ter sido a mais difícil para os estudantes.
Destaca -se, também, a questão 28, com Índice de Facilidade 0,14, o que, em termos percentuais, corresponde a 14 \% de estudantes que responderam acertadamente, com 0,01 de Índice de Discriminação, considerado Fraco.
Tais questões foram, portanto, pelo critério Ponto-Bisserial, consideradas inadequadas.
Por isso, as questões 10 e 28, além das questões 9, 12, 13, 22, 27, 31 e 34, foram eliminadas do cômputo da nota final.
A título de exemplo das análises do comportamento das questões objetivas, no Gráfico 6.5, analisa -se a questão 32 do Componente de Conhecimento Específico.
Essa questão foi considerada, pelas respostas dos estudantes avaliados na prova, como uma questão com nível de facilidade Médio, 0,44, ou seja, 44 \% dos estudantes assinalaram, acertadamente, a opção D, correspondente ao gabarito.
Como já comentado, seu Índice de Discriminação foi igual a 0,42, classificado como Muito bom, o maior valor de discriminação.
Neste gráfico, cada uma das cinco curvas representa o percentual de respostas em determinada alternativa da questão 32, em função do número de acertos dos estudantes nessa parte da prova, antes de possíveis eliminações de questões pelo critério Ponto-Bisserial.
A alternativa correta D, representada no gráfico pela curva em \textbf{laranja}, foi escolhida em maiores proporções pelos alunos com desempenho melhor nessa parte da prova.
Já as alternativas incorretas, também denominadas distratores, foram selecionadas, principalmente, por aqueles com notas mais baixas.
Observa -se que a soma não é 100 \%, por causa das questões não respondidas ou com mais de uma opção marcada.
Aqueles com nota zero, na sua quase totalidade, deixaram esta questão em branco ou marcaram mais de uma alternativa, comportamento considerado inválido.
A proporção de alunos que selecionaram a resposta correta D aumenta gradativamente, chegando a atingir 100 \% para 18 acertos ou mais, enquanto a proporção dos que escolheram alternativas incorretas decai, a partir de um ou dois acertos, em função do número de acertos nessa parte da prova.
Os gráficos relativos às demais questões do Conhecimento Específico constam do Anexo I. 6.3 ANÁLISE DAS QUESTÕES DISCURSIVAS Esta seção apresenta estatísticas selecionadas e histogramas das Questões Discursivas de Formação Geral ( 6.3.1 ) e Conhecimento Específico ( 6.3.2 ).
São, também, apresentadas e comparadas as médias de subpopulações, caracterizadas por Grande Região. 6.3.1 Componente de Formação Geral As análises dos resultados de desempenho dos estudantes de Ciências Contábeis nas duas questões discursivas relativas à Formação Geral encontram -se na Tabela 6.13 e no Gráfico 6.6.
Na Tabela 6.13, observa -se que a nota Média nesse conjunto de questões foi abaixo da obtida nas objetivas.
Os estudantes de todo Brasil obtiveram, em Formação Geral, Média 45,9 nas questões objetivas e 34,5, nas questões discursivas.
Pode -se notar, também, que o Desvio padrão nesse conjunto de questões foi menor do que o obtido nas objetivas : 22,1, nas questões objetivas e 19,0, nas questões discursivas.
A maior Média foi obtida na região Sul ( 35,4 ), e a menor, na região Centro-Oeste ( 32,2 ).
A Mediana de todo o Brasil, neste componente, foi 37,0, a mesma sendo obtida nas regiões Nordeste e Sudeste.
Na região Centro-Oeste, a Mediana foi menor ( 35,0 ), e, nas demais regiões, foi maior que o valor para o Brasil : Norte e Sul ( 38,0 ).
A nota Máxima ( 95,0 ) foi obtida na região Sudeste, sendo 83,5 a Máxima na região Norte ; 86,5, na Nordeste ; 88,5, na Sul ; e 94,5, na região Centro-Oeste.
A nota Mínima ( 0,0 ) foi a mesma em todas as regiões do Brasil.
No Gráfico 6.6, está representada a distribuição das notas nas questões discursivas do Componente de Formação Geral.
A moda dessa distribuição ocorre no intervalo ( 40 ; 50 ], seguida pelo intervalo ( 30 ; 40 ].
Destaca -se, também, como um máximo local, o intervalo [ 0 ; 10 ], com distribuição em torno de 15 \% do total de notas, sendo que, no intervalo [ 0 ; 10 ], inclui -se, além da nota zero, a frequência de alunos que deixaram esse tipo de questão em branco.
Na sequência, os resultados verificados, para cada uma das questões discursivas de Formação Geral, serão apresentados, estabelecendo -se relações com os conteúdos abordados em cada uma delas.
Os comentários da Banca de docentes corretores, a respeito do observado na correção das respostas dos estudantes, suas impressões e conclusões serão apresentados junto à análise de cada questão.
Cumpre esclarecer que, tendo -se em vista que as questões discursivas de Formação Geral são padronizadas, ou seja, constam de todas as provas, os comentários da Banca são os mesmos para todas as carreiras acadêmicas, sendo direcionados a todos os estudantes que participaram do Enade/2018.
A seguir, serão analisados os desempenhos da Área de Ciências Contábeis nas duas questões discursivas de Formação Geral do Enade/2018, comparando -se os resultados obtidos com comentários para cada questão. 6.3.1.1 Análise de Conteúdo da Questão Discursiva 1 do Componente de Formação Geral Os dados de Ciências Contábeis, obtidos a partir das respostas à questão 1, encontram -se na Tabela 6.14 e no Gráfico 6.7.
Nessa questão – de desempenho muito semelhante ao da discursiva 2 de Formação Geral – os alunos de todo o Brasil tiveram Média 30,2.
A maior Média para a questão 1 foi obtida na região Norte ( 31,3 ), e a menor, na região Centro-Oeste ( 28,3 ).
Quanto à variabilidade das notas, o Desvio padrão de todo o Brasil foi 22,5.
O menor Desvio padrão foi obtido na região Sul ( 21,9 ), e o maior Desvio padrão foi obtido na região Nordeste ( 23,6 ).
A Mediana, notas Máxima e Mínima da questão discursiva 1 foram iguais para todas as regiões do Brasil, 30,0, 100,0 e zero, respectivamente.
No Gráfico 6.7, mostra -se a distribuição das notas na questão discursiva 1 do Componente de Formação Geral.
A maior frequência ocorre no intervalo ( 40 ; 50 ], sendo este o modal.
Vale a pena ressaltar que o intervalo imediatamente anterior, ( 30 ; 40 ], tem uma frequência muito próxima.
Observa -se, também, que a frequência correspondente a dos alunos que deixaram a resposta a essa questão em branco foi elevada, 14,5 \% dos participantes, caracterizando -se como um máximo local. 6.3.1.2 Comentários sobre a correção de Conteúdo das respostas à Questão Discursiva 1 A Questão 1 exigia do aluno um posicionamento sobre conteúdo humanista e crítico.
O comando da questão solicitava a redação de um texto que formulasse e articulasse argumentos consistentes envolvendo temas previstos na Portaria nº 444, de 30 de maio de 2018, que dispõe sobre o componente de Formação geral do Enade 2018, tais como : solidariedade/violência, tolerância/intolerância, inclusão/exclusão, sexualidade, relações de gênero e relações étnico-raciais ( art. 7º, item VIII ), bem como ética, democracia e cidadania ( art. 7º, item I ).
O enunciado da questão buscava o estabelecimento da relação de três dimensões – os Direitos Humanos, a Democracia e as Fake News.
Para isso, recorreu a quatro suportes : a reprodução de uma foto e três fragmentos adaptados de textos que foram disponibilizados em páginas da Internet ( da Anistia Internacional, do Jornal El País-Brasil e do Labic-Laboratório de estudos sobre Imagem e Cibercultura ).
Os fragmentos traziam dados sobre o aumento do número de assassinatos de grupos minoritários, as origens e o trabalho da vereadora Marielle Franco, sua luta pelos Direitos Humanos, seu assassinato, e a produção de fake news.
O comando solicitava um texto abordando dois tópicos : o tensionamento entre a defesa dos Direitos Humanos e a produção de notícias falsas, tendo como base a vereadora Marielle – a defesa dos Direitos Humanos realizada por ela e as notícias falsas divulgadas após a sua morte.
As fake news também estavam presentes no segundo tópico de análise : os prejuízos que provocam para a sociedade democrática.
Por um lado, os textos selecionados para constituir o enunciado da questão continham apoio suficiente para formulação e articulação de argumentos para discorrer sobre o tensionamento entre a atuação da Marielle em defesa dos Direitos Humanos e fake news.
Por outro, havia ênfase nos aspectos biográficos e na trajetória da vereadora Marielle, o que levou a respostas pouco focadas no comando, que discutiam a vida e a morte da vereadora, sua importância como um ícone dos Direitos Humanos, opiniões em torno do seu assassinato, até mesmo sobre a responsabilidade pelo crime, muitos textos com a presença/defesa de fake news ou produção de novas.
O foco sobre Marielle também ensejou respostas basicamente centradas na questão da violência.
Foram frequentes as interpretações com viés político-partidário, não previstas no padrão de resposta.
O item do comando que solicitava a discussão sobre democracia e fake news encontrou menos apoio nos textos.
Apesar de se esperar que graduandos tivessem maturidade, inclusive política, para demonstrar alguma compreensão dos aspectos principais que envolvem o convívio democrático, esse aspecto foi pouco discutido e, quando presente, quase sempre foi apenas mencionado, numa quase reprodução do comando.
A falta de entendimento do que é o conceito de Direitos Humanos também foi percebida nas respostas, seja para recomendá -lo como importante para a sociedade, seja para criticá -lo como algo que defende bandidos.
Foi marcante a ideia que vincula Direitos Humanos a direitos apenas das minorias.
Não foi incomum encontrar respostas que personificavam os Direitos Humanos, tratando -os como sujeito de alguma ação que deveria ser implementada na sociedade em geral ou em casos específicos.
As fake news foram citadas muitas vezes, mas pouco explicadas, conforme previsto no padrão de resposta.
O argumento utilizado geralmente vinculava as consequências de fake news para a imagem da Marielle, ou seja, centrava no fato em si, já indicado no enunciado, sem qualquer análise ou aprofundamento próprio.
É importante salientar que muitas respostas, provavelmente de grupos vinculados a cursos específicos, continham argumentações recorrentes : os Direitos Humanos como algo vinculado à Constituição e a tratados dos quais o Brasil é signatário ; as fake news como vereadora Marielle Franco, sua luta pelos Direitos Humanos, seu assassinato, e a produção de fake news.
O comando solicitava um texto abordando dois tópicos : o tensionamento entre a defesa dos Direitos Humanos e a produção de notícias falsas, tendo como base a vereadora Marielle – a defesa dos Direitos Humanos realizada por ela e as notícias falsas divulgadas após a sua morte.
As fake news também estavam presentes no segundo tópico de análise : os prejuízos que provocam para a sociedade democrática.
Por um lado, os textos selecionados para constituir o enunciado da questão continham apoio suficiente para formulação e articulação de argumentos para discorrer sobre o tensionamento entre a atuação da Marielle em defesa dos Direitos Humanos e fake news.
Por outro, havia ênfase nos aspectos biográficos e na trajetória da vereadora Marielle, o que levou a respostas pouco focadas no comando, que discutiam a vida e a morte da vereadora, sua importância como um ícone dos Direitos Humanos, opiniões em torno do seu assassinato, até mesmo sobre a responsabilidade pelo crime, muitos textos com a presença/defesa de fake news ou produção de novas.
O foco sobre Marielle também ensejou respostas basicamente centradas na questão da violência.
Foram frequentes as interpretações com viés político-partidário, não previstas no padrão de resposta.
O item do comando que solicitava a discussão sobre democracia e fake news encontrou menos apoio nos textos.
Apesar de se esperar que graduandos tivessem maturidade, inclusive política, para demonstrar alguma compreensão dos aspectos principais que envolvem o convívio democrático, esse aspecto foi pouco discutido e, quando presente, quase sempre foi apenas mencionado, numa quase reprodução do comando.
A falta de entendimento do que é o conceito de Direitos Humanos também foi percebida nas respostas, seja para recomendá -lo como importante para a sociedade, seja para criticá -lo como algo que defende bandidos.
Foi marcante a ideia que vincula Direitos Humanos a direitos apenas das minorias.
Não foi incomum encontrar respostas que personificavam os Direitos Humanos, tratando -os como sujeito de alguma ação que deveria ser implementada na sociedade em geral ou em casos específicos.
As fake news foram citadas muitas vezes, mas pouco explicadas, conforme previsto no padrão de resposta.
O argumento utilizado geralmente vinculava as consequências de fake news para a imagem da Marielle, ou seja, centrava no fato em si, já indicado no enunciado, sem qualquer análise ou aprofundamento próprio.
É importante salientar que muitas respostas, provavelmente de grupos vinculados a cursos específicos, continham argumentações recorrentes : os Direitos Humanos como algo vinculado à Constituição e a tratados dos quais o Brasil é signatário ; as fake news como produto do mundo digital e globalizado, relacionado à era da pós-verdade ; a discussão sobre o papel das mídias sociais e da imprensa. produto do mundo digital e globalizado, relacionado à era da pós-verdade ; a discussão sobre o papel das mídias sociais e da imprensa.
Como já apontado, houve baixa quantidade de respostas contemplando o segundo tópico do comando – prejuízos da produção de notícias falsas para a sociedade democrática ; fato motivado, possivelmente, pela falta de um elemento que ressaltasse esse ponto nos textos motivadores.
De modo geral, os estudantes não conseguiram recorrer a conhecimentos tácitos sobre democracia para construir argumentos próprios.
As respostas que abordavam o tema, geralmente estavam vinculadas à ideia do direito à verdade, do direito à integridade da pessoa ou dos direitos gerais garantidos pela Constituição.
A relação da questão com o momento político brasileiro foi indicada em um número expressivo de respostas.
Por vezes, apenas questionando a temática da questão – caso Marielle – em uma prova para avaliar o ensino superior no Brasil ou a importância dos Direitos Humanos na atual conjuntura política do país, levando à anulação da resposta.
Houve também interpretações das consequências do uso das fake news no processo eleitoral de 2018, influenciando a escolha consciente e livre das pessoas.
As respostas demonstraram atenção político social frente a esses temas, porém com pouco domínio conceitual do campo de conhecimento sobre Direitos Humanos, fake news e democracia.
Em termos da linguagem adotada pelos alunos, há dois grandes grupos de respostas.
No primeiro grupo, os que responderam de forma clara, com autonomia de pensamento e explicitação de justificativas para os aspectos solicitados no enunciado, demonstrando conhecimento sobre o que significam Direitos Humanos, fake news e democracia, além da articulação entre essas questões.
Por outro lado, um segundo grupo dos que não conseguiram ir além do que foi citado no enunciado, apenas repetindo -o, muitas vezes desviando o foco da questão para temas correlatos, como a vida de Marielle ou a violência.
Nesse segundo grupo, mais facilmente foram encontradas respostas desconexas e mal articuladas, com textos confusos que explicitam o não conhecimento sobre o tema e a falta de adequação à linguagem, necessária a um concluinte do ensino superior.
As abordagens dos Direitos Humanos mais recorrentes nas respostas corretas foram : a referência da Marielle como representante/ícone dos DH ; e os DH como direitos previstos na Constituição e em tratados internacionais dos quais o Brasil é signatário.
No entanto, também foram encontradas críticas aos Direitos Humanos, tais como : DH vistos como “ direitos de bandidos ” ; a defesa pouco embasada de frases de efeito como “ direitos humanos para humanos direitos ”.
Também chamou atenção o tratamento dos DH como um personagem, ou um órgão público, evidenciando desconhecimento sobre o tema.
Em relação às fake news, dois argumentos previstos pelo padrão de respostas foram os mais presentes : a noção de que o mundo globalizado e a facilidade de reprodução de informações pelas mídias atuais levam à criação de mentiras que são compartilhadas e que geram consequências ; a Marielle como alvo de fake news que tentaram desconstruir sua luta e seu legado.
Também foi frequente a ideia de que a produção de informações falsas é realizada por uma imprensa tendenciosa com intenção de manipular a informação.
No item relacionado a fake news e democracia, as respostas corretas afirmavam que as fake news inviabilizam o acesso à verdade e à garantia da integridade de um Estado Democrático de Direito.
Também foi abordada a influência das fake news nos resultados das eleições dos últimos anos, nos Estados Unidos e no Brasil, gerando consequências para a vida das pessoas.
Por fim, um outro aspecto abordado nas respostas a este tópico foi a polarização da sociedade em termos políticos, com a afirmação de que as fake news dificultam o debate, estimulam posições extremas e rivalidades prejudiciais à convivência em uma democracia plena.
A notas mais fracas ( de zero a 35 ) foram cerca de 20 \% das respostas corrigidas – as que não estavam em branco nem foram desconsideradas ou anuladas.
Foram respostas que se restringiam a cópia ou paráfrases do enunciado ; as que tinham foco em desdobramentos políticos e policiais do caso Marielle ; as que tinham argumentação confusa e/ou incompleta.
A grande maioria das notas, em torno de 65 \%, foi mediana, entre 40 a 70.
Nesse grupo ficaram as respostas que continham alguma abordagem adequada sobre Direitos Humanos, sobre as atividades da vereadora Marielle e sobre as fake news divulgadas após sua morte.
As notas mais altas deste grupo foram para respostas que incluíram a questão do tensionamento, quase sempre, em um texto coerente e consistente.
Muitas vezes, a nota mediana foi consequência do baixo índice de respostas à segunda parte do comando, que solicitava a discussão sobre os prejuízos das fake news para a democracia.
As melhores notas só foram obtidas por aqueles que incluíram a discussão sobre os prejuízos das fake news para a democracia.
Apenas 15 \% das respostas corrigidas receberam avaliação igual ou superior a 75.
A deficiência principal detectada pelos corretores foi a dificuldade de grande parte dos alunos para interpretar o enunciado da questão e para se expressar em textos consistentes.
De modo geral, os temas foram tratados de forma fragmentada e pouco analítica.
Além disso, chamou à atenção a predominância do lugar comum na construção argumentativa das respostas, a superficialidade do debate em torno da democracia e o desconhecimento de conceitos como Direitos Humanos, esses confundidos com instituições e até pessoas. 6.3.1.3 Análise de Conteúdo da Questão Discursiva 2 do Componente de Formação Geral A Tabela 6.15 mostra que o desempenho médio dos estudantes na questão discursiva 2 ( média 30,5 ) foi muito próximo ao obtido na questão discursiva 1 ( média 30,2 ).
A região Sul foi aquela cuja Média nessa questão foi maior ( 32,1 ), enquanto a de menor Média, também, foi a da região Centro-Oeste ( 28,3 ), inclusive com a mesma Média obtida na discursiva 1.
Quanto à variabilidade das notas, o Desvio padrão de todo o Brasil foi 22,5, igual ao obtido na questão discursiva 1.
O maior desvio, por sua vez, nessa questão, foi obtido na região Norte ( 23,2 ), enquanto o menor foi obtido na região Sul ( 22,2 ).
A Mediana de todo o Brasil foi 30,0, a mesma em quase todas as regiões.
A exceção foi a região Sul, que obteve Mediana 35,0.
As notas : Máxima ( 100,0 ) e Mínima ( 0,0 ) foram as mesmas em todas as regiões, sem exceção.
No Gráfico 6.8, mostra -se a distribuição das notas na questão discursiva 2 do Componente de Formação Geral.
Observa -se que a maior frequência corresponde ao intervalo de notas ( 30 ; 40 ], seguido do intervalo ( 40 ; 50 ].
A frequência de alunos que deixaram essa questão em branco ultrapassa os 16 \% e figura na distribuição como um máximo local.
No tocante à dispersão, em comparação à questão discursiva 1, o comportamento é similar, pois ambas as questões têm o mesmo Desvio padrão ( 22,5 ). 6.3.1.4 Comentários sobre a correção de Conteúdo das respostas à Questão Discursiva 2 Os textos e o comando da Questão 2 de Formação Geral são perfeitamente adequados ao que se espera avaliar em Formação Geral, ao tratar das funções sociais dos museus, contextualizando a temática sob o ponto de vista da memória e da perspectiva de futuro.
Na questão, estavam presentes, principalmente, o conteúdo humanista e crítico, a interpretação e a produção de textos sobre temas como cultura e arte, educação, ciência e democracia, previstos na Portaria que dispõe sobre o componente de Formação Geral do Enade 2018.
Os conhecimentos exigidos para responder ao comando, explicitados no padrão de resposta, vão além do senso comum, por se exigir a compreensão de conceitos como “ função social ” ( de instituições sociais ), “ memória ” ( no sentido de memória social ), “ perspectiva de futuro ” ( no contexto sócio-histórico-cultural ) e “ sociedade contemporânea ” ( que deveria independer da interpretação de cada grupo social a respeito da sociedade na qual vive ).
No entanto, apesar de envolver tais conceitos, espera -se que alunos concluintes de cursos superiores de graduação sejam capazes de discutir funções de instituições sociais, especialmente porque o enunciado trazia fato recente e de grande repercussão como apoio.
O padrão de resposta elencou as funções sociais que poderiam ser abordadas nas respostas dos estudantes, como também argumentos que poderiam ser usados ao longo das mesmas.
A tarefa dos estudantes era a construção de um texto expositivo que articulasse a importância de museus, suas funções sociais e as perspectivas de memória e de futuro, o que se mostrou de dificuldade de média a difícil.
Apesar de o enunciado ser claro, muitos estudantes, provavelmente influenciados pelos textos, limitaram -se a discutir aspectos relacionados ao Museu Nacional e, em particular, ao incêndio que o destruiu.
Foram poucos os que interpretaram corretamente o escopo do comando, que solicitava a discussão das funções sociais de museus de uma forma geral, no sentido de argumentar a importância deste tipo de instituição.
As quatro funções sociais elencadas no padrão de resposta foram encontradas na quase totalidade das redações dos estudantes de forma articulada com a importância dos museus.
Entretanto, quase nunca os textos abordavam mais do que uma função social e vinculavam, de forma clara, a importância desse tipo de instituição para a sociedade contemporânea.
Além disso, a solicitação de contextualização sob o foco museu/memória/perspectiva de futuro parece não ter sido bem compreendida e, em consequência, esteve pouco presente nas respostas.
Também se observou confusão entre os conceitos de história e memória e uma visão estreita do significado de sociedade contemporânea, às vezes, restrita à cidade do respondente, outras, ao Rio de Janeiro, ou ainda a um tipo de sociedade idealizada.
A relevância que os textos de apoio deram ao caso particular do incêndio do Museu Nacional do Rio de Janeiro e aos prejuízos causados ao seu acervo, induziu a maioria dos estudantes a elaborar respostas tratando das funções a partir da perda e, em particular, à perda do tipo de acervo daquele museu e do tipo de pesquisa que lá se desenvolve, divergindo do comando da questão, que solicitava texto abordando funções sociais de museus em geral.
Também foram encontradas muitas respostas que se limitavam a discutir o que o estudante considerava como possíveis causas do incêndio, a falta de manutenção, questões políticas e de gestão, sem responder ao que foi solicitado.
Foi também recorrente a transcrição de frases dos textos do enunciado nas respostas, diretamente ou por paráfrases, revelando certa dificuldade de criação textual.
A análise das respostas dos estudantes, além de revelar dificuldade de abordar a temática esperada, revelou outros dois aspectos predominantes.
Um deles, como já comentado, foi a ênfase no incêndio do Museu Nacional, que orientou grande parte dos textos superiores de graduação sejam capazes de discutir funções de instituições sociais, especialmente porque o enunciado trazia fato recente e de grande repercussão como apoio.
O padrão de resposta elencou as funções sociais que poderiam ser abordadas nas respostas dos estudantes, como também argumentos que poderiam ser usados ao longo das mesmas.
A tarefa dos estudantes era a construção de um texto expositivo que articulasse a importância de museus, suas funções sociais e as perspectivas de memória e de futuro, o que se mostrou de dificuldade de média a difícil.
Apesar de o enunciado ser claro, muitos estudantes, provavelmente influenciados pelos textos, limitaram -se a discutir aspectos relacionados ao Museu Nacional e, em particular, ao incêndio que o destruiu.
Foram poucos os que interpretaram corretamente o escopo do comando, que solicitava a discussão das funções sociais de museus de uma forma geral, no sentido de argumentar a importância deste tipo de instituição.
As quatro funções sociais elencadas no padrão de resposta foram encontradas na quase totalidade das redações dos estudantes de forma articulada com a importância dos museus.
Entretanto, quase nunca os textos abordavam mais do que uma função social e vinculavam, de forma clara, a importância desse tipo de instituição para a sociedade contemporânea.
Além disso, a solicitação de contextualização sob o foco museu/memória/perspectiva de futuro parece não ter sido bem compreendida e, em consequência, esteve pouco presente nas respostas.
Também se observou confusão entre os conceitos de história e memória e uma visão estreita do significado de sociedade contemporânea, às vezes, restrita à cidade do respondente, outras, ao Rio de Janeiro, ou ainda a um tipo de sociedade idealizada.
A relevância que os textos de apoio deram ao caso particular do incêndio do Museu Nacional do Rio de Janeiro e aos prejuízos causados ao seu acervo, induziu a maioria dos estudantes a elaborar respostas tratando das funções a partir da perda e, em particular, à perda do tipo de acervo daquele museu e do tipo de pesquisa que lá se desenvolve, divergindo do comando da questão, que solicitava texto abordando funções sociais de museus em geral.
Também foram encontradas muitas respostas que se limitavam a discutir o que o estudante considerava como possíveis causas do incêndio, a falta de manutenção, questões políticas e de gestão, sem responder ao que foi solicitado.
Foi também recorrente a transcrição de frases dos textos do enunciado nas respostas, diretamente ou por paráfrases, revelando certa dificuldade de criação textual.
A análise das respostas dos estudantes, além de revelar dificuldade de abordar a temática esperada, revelou outros dois aspectos predominantes.
Um deles, como já comentado, foi a ênfase no incêndio do Museu Nacional, que orientou grande parte dos textos dos estudantes.
Outro, também evidenciando desvio do esperado, foi o foco na perda patrimonial.
Nesse grupo, houve respostas que até apresentaram uma função social, mas olhada pelo ponto de vista da perda.
Nessa linha, muitas respostas enalteciam a importância do Museu Nacional, sua relevância no cenário nacional, o impacto da perda de seu acervo, isso complementado por frágeis reflexões sobre as funções sociais.
Em consequência, a perspectiva de desdobramentos mais ricos, que tratassem de aspectos sociais relacionados a diferentes tipos de museu foi prejudicada.
Muitas respostas trataram da preservação da memória e/ou guarda do acervo ; identidade, tanto individual quanto coletiva ; conhecimento sobre antepassados ; evolução da sociedade e dos seus costumes ; valorização e disseminação do conhecimento por meio da liberação do acesso aos acervos ; bem como atividades pedagógicas e exposições de arte.
O aspecto produção de pesquisa e conhecimento, embora frequente, foi mencionado de forma muito simples, sem qualquer aprofundamento, para além do que já era mencionado no enunciado.
Frequentemente apresentava -se o museu como produtor de pesquisas para o campo da saúde, o que evidencia desconhecimento das reais possibilidades de pesquisa, mesmo no Museu Nacional.
Quanto a perspectiva de futuro, quando presente no texto, era mencionada de forma breve, com afirmações genéricas sobre a possibilidade de melhoria do futuro pelo simples conhecimento do passado, numa percepção de que o passado determina o futuro, de forma linear.
Foram raras as respostas com alguma argumentação que associasse pesquisa, memória, identidade e reflexão social como base para a construção de alternativas de futuro.
A relação entre turismo e museu apareceu considerando a visita a museus como forma de lazer e, algumas vezes, como geradora de recursos.
O desempenho pode ser considerado fraco.
Muitas respostas deixavam claro que o estudante nunca tinha visitado um museu e, de uma forma geral, observou -se desconhecimento sobre suas funções sociais.
Não se esperava que história, memória social e museologia fossem assuntos dominados pelos estudantes das diversas áreas que participaram do Enade 2018, mas esperava -se que essa temática pudesse ser abordada do ponto de vista das experiências pessoais ou de conhecimento social.
Como uma vivência própria com os museus também se revelou frágil ou inexistente, o conhecimento sobre o assunto ficou bem restrito aos textos motivadores e ao senso comum, construído, basicamente, pelo noticiário em torno do incêndio.
O universo dos argumentos dos estudantes se mostrou muito limitado.
Cada dimensão foi praticamente representada por uma ideia simplista.
De forma recorrente : a preservação da memória para construção da identidade de um povo ; a compreensão do passado por meio das pesquisas como oportunidade para melhorar o futuro ; o museu como um espaço em que há um aprendizado estimulado pelo lazer.
Houve, também, uma grande concentração de respostas associando o museu a um local de guarda de acervo e de história.
As respostas à Questão 2 evidenciaram dificuldades conceituais.
Além dos conceitos de história, memória e museu se misturarem, parecendo não se distinguirem, ao mencionarem a pesquisa, também se observou falta de distinção entre pesquisas colegiais, com objetivos de aprendizado, e as realizadas por pesquisadores, com objetivos de avançar cientificamente.
A falta de compreensão do que se pedia no comando da questão e a dificuldade de expressão linguística foram dois aspectos marcantes na avaliação dos corretores.
Foi frequente os textos apresentarem ideias soltas ou de forma confusa, sem constituir unidade.
O vocabulário usado se revelou bem restrito, e fortemente apoiado no coloquial, na linguagem oral.
Predominou o senso comum em lugar de uma reflexão mais formal envolvendo os conceitos envolvidos. 6.3.1.5 Análise de Língua Portuguesa das Questões Discursivas do Componente de Formação Geral Os dados de Ciências Contábeis, obtidos a partir das respostas às questões discursivas do Componente de Formação Geral, no que tange à Língua Portuguesa, encontram -se na Tabela 6.16 e no Gráfico 6.9.
Nesse aspecto, os alunos de todo o Brasil obtiveram Média 51,5.
A maior Média em relação à Língua Portuguesa foi obtida na região Sudeste ( 52,6 ), e a menor, na região Centro-Oeste ( 48,0 ).
Quanto à variabilidade das notas, o Desvio padrão de todo o Brasil foi 24,0.
O menor Desvio padrão foi obtido na região Sul ( 22,8 ) e o maior Desvio padrão foi obtido na região Centro-Oeste ( 25,1 ).
A Mediana das notas de Língua Portuguesa foi 60,0 para o Brasil, como um todo, a mesma obtida nas regiões Norte, Sudeste e Sul.
A Mediana das regiões Nordeste e Centro-Oeste foi 57,5.
A nota Máxima para todo o Brasil foi 95,0, com, pelo menos, um aluno tirando essa nota na região Nordeste.
Na região Norte, a nota Máxima foi 87,5 ; e nas regiões Sudeste, Sul e Centro-Oeste foi 92,5.
Já a nota Mínima foi zero em todas as regiões do país.
No Gráfico 6.9, mostra -se a distribuição das notas de Língua Portuguesa do Componente de Formação Geral.
Observa -se que a maior frequência ( acima de 28 \% ) corresponde à dos alunos que obtiveram nota no intervalo ( 60 ; 70 ].
Destacam -se, também, os alunos que deixaram ambas as questões em branco, representando um pouco mais do que 11 \% do total, caracterizando -se como máximo local. 6.3.1.6 Comentários sobre a correção das respostas de Formação Geral com respeito à Língua Portuguesa As duas questões discursivas de Formação Geral ensejaram duas configurações textuais distintas, tendo em vista os seus encaminhamentos : a Questão 1 teve um caráter mais opinativo e a Questão 2, mais expositivo, conforme se explicita a seguir.
Ambas as questões permitiram a observação do desempenho linguístico dos participantes em situação formal, competência exigida no processo de formação universitária.
As habilidades de produção textual envolvidas têm, igualmente, como condição básica, o domínio do padrão formal culto da Língua Portuguesa, no que diz respeito aos aspectos relativos às questões de caráter ortográfico, textual, morfossintático e vocabular.
O encaminhamento das questões evidencia a expectativa de que o participante utilizasse seus conhecimentos sobre os dois temas e estruturasse seus textos de acordo com as características do registro formal adequado à situação comunicativa.
Tal configuração determina exigências quanto : à adequação da seleção vocabular, ao desenvolvimento do conteúdo, à estruturação sintática dos períodos, à organização lógica das ideias, à utilização de procedimentos de encadeamento textual e de referenciação, à obediência às exigências morfossintáticas próprias da modalidade escrita da norma-padrão, ao respeito às regras ortográficas e às regras de acentuação gráfica.
Esse encaminhamento exigiu do graduando o domínio do modo de organização textual de base dissertativo-argumentativa, essencial ao processo de formação universitária.
As habilidades de produção textual envolvidas têm como condição básica o domínio do padrão formal culto da Língua Portuguesa, no que diz respeito aos aspectos relativos às questões de caráter ortográfico, textual, morfossintático e vocabular.
O enunciado da Questão 1 tomava como base três fragmentos de textos jornalísticos publicados em sites da internet, que abordavam o aumento do assassinato de minorias e de defensores de direitos humanos no Brasil, entre eles a vereadora Marielle Franco.
O comando solicitava que o participante desenvolvesse um texto dissertativo e argumentativo com as temáticas Fake News, Democracia e Direitos Humanos.
Embora não tenha havido referência explícita à tipologia textual esperada para o desenvolvimento da resposta, essa proposta encaminhou o participante para a elaboração de um texto opinativo, com o objetivo de discutir o problema e manifestar um posicionamento crítico.
O enunciado da Questão 2 abordava o incêndio do Museu Nacional do Rio de Janeiro, ocorrido em 2018, e seus impactos na ciência brasileira e internacional, tomando como base dois fragmentos de textos jornalísticos publicados em sites da internet.
O comando solicitava que o participante desenvolvesse um texto abordando três aspectos da função social dos museus na sociedade contemporânea, encaminhando -o para a elaboração de um texto expositivo que confirmasse a posição defendida nos fragmentos motivadores a respeito da importância dessas instituições.
Para atender a essas exigências, o padrão de resposta utilizado na avaliação das questões discursivas de Formação Geral considerou aspectos relevantes ao bom desempenho linguístico como competências distintas, de modo a permitir um mapeamento detalhado do domínio dos recursos disponíveis na Língua Portuguesa para a comunicação escrita formal : ortográficos, textuais, morfossintáticos e vocabulares ( ver detalhes no Anexo VIII – Padrão de Respostas ).
Aspectos Ortográficos Para avaliar a competência relativa ao domínio das convenções ortográficas da norma-padrão da Língua Portuguesa, observou -se o respeito às regras de acentuação gráfica e da grafia padrão das palavras ( com ausência de abreviaturas próprias da linguagem da internet ), de acordo com as convenções estabelecidas pela legislação em vigor e consubstanciadas no Vocabulário Ortográfico da Língua Portuguesa, editado pela Academia Brasileira de Letras.
Aspectos Textuais Para avaliar a competência relativa ao domínio dos procedimentos de estruturação textual, procurou -se observar a utilização de mecanismos linguístico-discursivos que estabeleçam o encadeamento lógico entre as partes do texto, de forma a garantir a progressão e a coerência textuais.
Esta competência destaca -se por sua relevância, já que são avaliadas a clareza e a unidade textual.
Aspectos Morfossintáticos Para avaliar o domínio dos diferentes aspectos morfossintáticos próprios da modalidade escrita formal da norma-padrão da Língua Portuguesa, observou -se se o texto produzido atendeu às seguintes exigências : a concordância nominal, a concordância verbal, a regência nominal, a regência verbal, a flexão nominal, a flexão verbal, a correlação entre os tempos verbais, a colocação pronominal e a utilização de sinais de pontuação que contribuam para a organização lógica da frase e do texto.
Aspectos Vocabulares Para avaliar a adequação da seleção vocabular à modalidade escrita formal da Língua Portuguesa, exigida pela situação comunicativa, observou -se o atendimento às seguintes exigências : precisão na seleção/utilização do vocabulário relacionado à temática solicitada pela questão ; ausência de marcas da oralidade, como termos de sentido muito genérico ( “ coisa ”, “ negócio ”, “ você ” ) e termos de registros mais informais ( como gírias, jargões, frases feitas, ditados populares, termos regionais ).
Assim, espera -se que o participante respeite a adequação vocabular não usando gírias ou expressões coloquiais, evite repetição desnecessária de palavras e utilize um vocabulário mais formal, como solicitado por um texto dissertativo.
Para efeito de pontuação, a avaliação do desempenho linguístico considerou três grandes grupos de competências, agrupando aspectos explicitados, com os seguintes pesos relativos : aspectos ortográficos ( 20 \% ) ; aspectos textuais ( 40 \% ), aspectos morfossintáticos e vocabulares ( 40 \% ).
Os critérios de pontuação utilizados no padrão de resposta procuraram atender à diversidade de desempenho dos graduandos.
Nos textos analisados, observaram -se, assim, vários níveis de atendimento às exigências do padrão formal próprio da modalidade escrita da Língua Portuguesa, configurando um"continuum".
Por um lado, textos que traduzem com exatidão a expectativa de uso formal da língua, em todos os níveis analisados, marcados por clareza e unidade ; e, no outro extremo, textos quase incompreensíveis devido ao grande número de desvios de todos os tipos, principalmente os de caráter textual.
As duas questões discursivas de Formação Geral analisadas revelaram desempenhos distintos dos participantes.
A Questão 1 propiciou a oportunidade de maior desenvolvimento do tema solicitado, gerando consequências na estruturação textual, já que os textos foram mais longos, mais elaborados e fluentes.
Em virtude da polêmica sobre a existência das fake news e sobre a visão pública da personagem da vereadora Marielle Franco, observaram -se respostas que polemizaram o tema, favorecendo a construção argumentativa do texto.
A Questão 2 teve um comportamento distinto, devido ao caráter mais expositivo do enunciado e ao reduzido conhecimento dos participantes sobre o tema.
Assim, os textos foram mais curtos, frequentemente com repetição de ideias dos textos motivadores.
Quanto aos aspectos linguísticos analisados durante esta avaliação, observaram -se os seguintes resultados : Aspectos ortográficos : O desempenho dos participantes revelou uma diferença muito grande nos dois aspectos analisados nesta competência : baixo índice de desvios da grafia padrão e grande índice de desvios de acentuação.
Em vários casos, ocorre ausência completa de acentuação gráfica em todas as palavras do texto.
Os resultados revelaram que a tendência dominante entre os universitários brasileiros é a eliminação da acentuação gráfica, talvez motivada pelos hábitos relacionados às redes sociais e pela ausência de esclarecimento dos meios de comunicação, das autoridades e das escolas sobre as decisões do Acordo Ortográfico da Língua Portuguesa de 1990.
Os casos mais sistemáticos de eliminação do acento indicador da sílaba tônica foram : - palavras proparoxítonas ( por exemplo “ generos ”, “ arbitro ”, “ politico ”, “ publicas ”, “ numero ” ) ; - palavras paroxítonas terminadas em ditongo crescente ( por exemplo “ area ”, “ noticiarios ”, “ individuo ”, “ historia ”, “ varios ”, “ noticia ”, “ negligencia ”, “ ocorrencia ” ) ; - palavras oxítonas ( por exemplo “ ninguem ”, “ esta ”, “ ate ”, “ tambem ” ).
Por outro lado, destaca -se o uso indevido do acento gráfico em palavras como “ melâncolia ” e “ intervênção ”, por exemplo.
Quanto ao domínio das convenções relativas à grafia das palavras, observaram -se desvios esporádicos como : hipercorreção pela escolha de “ e ” no lugar de “ i ”, por influência de hábitos da oralidade ( “ descriminação ” no lugar de “ discriminação ”, “ entervenção ” no lugar de “ intervenção ” ) ; desvios de grafia relacionados à variação diastrática, como por exemplo “ perca ”, “ subjulgadas ”, “ precoseito ”, “ soubre ”, “ vecendo ”, “ indesencia ”, “ apolojia ”, “ fulga ”, “ dereitos ” ;...
Observaram -se, também, muitos casos de inadequação no uso da maiúscula : uso indevido para destacar determinadas palavras-chave do texto, como “ Brasileiros ”, “ Homens ” ; ausência de maiúscula para grafar nomes próprios, como"marielle".
Expressões como"Direitos Humanos"foram grafadas diversificadamente, com maiúsculas e com minúsculas.
Destaque -se, também, a maioria dos participantes que grafam os textos inteiramente em caixa alta, o que impede a identificação de desvios relacionados ao uso de maiúsculas.
Vale observar, também, que, ao contrário do que se esperava, não apareceram abreviaturas próprias do “ internetês ”, relacionadas ao uso de redes sociais e e-mails, como por exemplo"vc".
Aspectos textuais : Para alguns dos graduandos essa competência se revelou como a mais problemática, tendo em vista os inúmeros problemas observados, desvios acumulados durante toda a formação escolar.
São eles : sequência justaposta de ideias sem encaixamentos sintáticos, provocando uma redução de estruturas subordinadas, ao lado do aumento na frequência de estruturas coordenadas e absolutas ; frases fragmentadas que comprometem a estrutura lógico-gramatical ; frases formadas apenas por oração subordinada, sem oração principal ; redução no uso de conectores para expressar relações lógicas essenciais à construção do texto como consequência da mudança de estruturação frasal ; emprego inadequado do pronome relativo ( com omissão da preposição ou a utilização do pronome relativo inadequado, como “ onde ” ), refletindo hábitos da oralidade ; ausência de recursos de referenciação, como substituição de termos por sinônimos, hiperônimos, nominalizações, expressões metafóricas. lógico-gramatical ; frases formadas apenas por oração subordinada, sem oração principal ; redução no uso de conectores para expressar relações lógicas essenciais à construção do texto como consequência da mudança de estruturação frasal ; emprego inadequado do pronome relativo ( com omissão da preposição ou a utilização do pronome relativo inadequado, como “ onde ” ), refletindo hábitos da oralidade ; ausência de recursos de referenciação, como substituição de termos por sinônimos, hiperônimos, nominalizações, expressões metafóricas.
Esses problemas revelam dificuldades em relação à estrutura formal do texto produzido, o que é preocupante ao se levar em conta que são graduandos em fase final de formação.
Em alguns textos, falta um mínimo de textualidade e de domínio do registro padrão da língua.
Na verdade, observaram -se relações linguísticas quase agramaticais, como as estabelecidas pela sequência de gerúndios sem o apoio de um ponto de partida para a organização das informações gramaticais e semânticas.
Quanto à utilização dos sinais de pontuação, observou -se uma ausência dos mesmos na maioria dos textos analisados.
Foi muito frequente a ocorrência de parágrafos sem marca interna de pontuação para separar os períodos.
São os seguintes os tipos de problemas mais encontrados :  vírgula : utilização de vírgula para separar o sujeito e o predicado ; uso de vírgula no lugar do ponto para separar ideias que constituem períodos distintos ; ausência de vírgula para separar oração adjetiva explicativa ;  ponto e vírgula : utilização do ponto e vírgula no lugar de vírgula ;  ponto final : ausência de ponto final para separar períodos.
Aspectos morfossintáticos e vocabulares : Os resultados são transparentes em relação aos aspectos mais problemáticos no desempenho dos participantes nestas duas competências.
Em relação aos aspectos morfossintáticos, seguem algumas observações sobre os desvios mais frequentes.
Quanto à regência, o desvio mais frequente foi a falta do sinal indicativo da crase – isso revela que o usuário não tem consciência de que, sob a forma do termo “ a ”, existe a presença de uma contração entre a preposição “ a ” ( exigida pela regência do termo anterior ) e o artigo definido “ a ”.
Outro problema relacionado à regência verbal e à nominal, encontrado frequentemente nas respostas, foi a ausência de preposição antes de pronome relativo, processo generalizado na modalidade oral da língua, em situações de registro informal.
Apesar da possibilidade de que essa alteração de regência se generalize no padrão escrito da Língua Portuguesa, como já está ocorrendo até em textos jornalísticos, o não emprego da preposição foi considerado como desvio neste processo de avaliação.
Outro desvio de regência significativo foi a utilização inadequada de uma preposição ou sua ausência após o verbo ou o nome ( substantivo ou adjetivo ).
A concordância verbal e a concordância nominal apresentaram alguns desvios muito frequentes.
Quanto à concordância de número, observou -se, como apontam as pesquisas nessa área, ausência de marca ( com sujeito anteposto ou posposto ) ou uso indevido ( uso inadequado da marca de plural comandado pelo núcleo plural da locução adjetiva, apesar de o substantivo que funciona como núcleo do sintagma nominal estar no singular ).
Uma ocorrência que se destacou foi a ausência de acento circunflexo na forma plural do presente do indicativo do verbo “ ter ”, que foi considerada como um desvio na concordância verbal e não na acentuação gráfica.
Quanto à concordância de gênero, vários casos foram observados, normalmente no âmbito de sintagmas nominais longos, em que o adjetivo está afastado do substantivo.
Quanto à questão da colocação pronominal, foram poucos os desvios observados.
Concluiu -se que, no registro escrito formal, a maioria dos participantes já incorporou regras como a não introdução da frase por um pronome oblíquo e a próclise na presença de um termo atrator.
Várias marcas de oralidade foram identificadas, embora não com alta frequência : uso do pronome relativo “ onde ” como relativo universal, falta de artigo definido antes de substantivo isolado ou antes de termos em enumerações.
Quanto aos aspectos vocabulares, alguns tipos de inadequação foram observados : expressões da oralidade ; seleção vocabular incompatível com o contexto, gerando falta de inteligibilidade ; falta de domínio de vocabulário mais abstrato e de maior complexidade, essencial ao desenvolvimento do texto de base dissertativa.
O principal aspecto observado foi a excessiva repetição de certas palavras, como o termo"pessoa", por exemplo, revelando limitação de repertório vocabular.
Em função do tema solicitado na Questão 1, diferentes grafias da expressão fake news não foram consideradas.
Para sintetizar, os problemas que mais se destacaram nos textos analisados foram : a ) truncamentos sintáticos, caracterizados por estruturas frasais incompreensíveis devido à ruptura da complexidade sintática própria da modalidade escrita ; b ) redução de estruturas subordinadas, compensada pelo aumento na frequência de estruturas coordenadas e absolutas, por um lado, ou pela elaboração de estruturas textuais prejudicadas pelo excesso de ideias sem a devida conexão subordinativa ; c ) redução no uso de conectores para expressar relações lógicas essenciais à construção do texto, substituídos pela exigência de inferência por parte do interlocutor para suprir a sua ausência ; d ) redução no uso do subjuntivo, ao lado da ampliação do uso do indicativo, combinado a estruturas frasais coordenadas ou absolutas ; d ) empobrecimento do processo de referenciação, com a repetição de palavras, sem processos mais sofisticados de substituição ; e ) simplificação da marcação da categoria tempo na morfologia verbal ; f ) desrespeito às regras de concordância verbal e nominal próprias do padrão formal ; g ) ausência do acento grave indicador da crase, revelando falta de conhecimento dos critérios de regência verbal e nominal ; h ) falta de domínio de vocabulário mais abstrato e de maior complexidade, essencial ao desenvolvimento do processo dissertativo ; i ) redução do emprego da acentuação gráfica.
Como resultado da correção, merece atenção a constatação de que a modalidade escrita tem apresentado a tendência a uma evidente simplificação, aproximando -se das características da modalidade oral da Língua Portuguesa.
No caso dos textos de base dissertativa ( expositiva ) e de base argumentativa, inscritos em um registro formal, a distância entre as duas modalidades é ainda maior, o que provoca desvios recorrentes em todos os quatro aspectos analisados : ortográficos, textuais, morfossintáticos e vocabulares.
Destaca -se, nessa avaliação, o comprometimento textual, por meio de estruturas fragmentadas e/ou truncadas, rompendo a complexidade sintática esperada no padrão formal. 6.3.2 Componente de Conhecimento Específico Na parte da prova relativa às questões discursivas do Componente de Conhecimento Específico ( Tabela 6.17 ), observa -se que a Média foi bem mais baixa do que para as questões discursivas do Componente de Formação Geral.
Enquanto no Componente de Formação Geral a Média para estudantes de Ciências Contábeis de todo o Brasil foi 34,5, na parte de Conhecimento Específico, a Média foi 14,6.
A maior Média deste componente foi obtida pelos estudantes da região Sudeste ( 16,8 ), e a menor, pelos da região Norte ( 10,6 ).
Quanto à Quanto à variabilidade das notas, o Desvio padrão de todo o Brasil foi 18,1.
O maior Desvio padrão foi encontrado na região Sudeste ( 19,4 ), e o menor, na região Norte ( 14,7 ).
A maior nota Máxima, 100,0, foi obtida em quase todas as regiões, com exceção da região Norte ( 95,0 ).
A nota Mínima ( 0,0 ) foi obtida por, pelo menos, um aluno em todas as regiões do Brasil.
A Mediana do Brasil, como um todo, foi 6,7.
A maior Mediana foi obtida na região Sudeste ( 10,0 ), e a menor, na região Norte ( 3,3 ).
No Gráfico 6.10, representa -se a distribuição das notas nas questões discursivas no Componente de Conhecimento Específico.
A moda dessa distribuição ocorre no intervalo [ 0 ; 10 ], com mais do que 56 \% do total de participantes.
Observa -se que nesse intervalo estão computadas as ocorrências de respostas em branco para as três questões discursivas de Conhecimentos Específicos.
A análise de cada uma destas questões será feita a seguir. 6.3.2.1 Análise da Questão Discursiva 3 do Componente de Conhecimento Específico Na questão 3, cujos resultados aferidos encontram -se descritos na Tabela 6.18, a Média dos estudantes de todo o Brasil foi 24,4, a questão de melhor desempenho entre as três discursivas de Conhecimentos Específicos.
A menor Média nessa questão foi obtida pelos alunos da região Norte ( 17,0 ), enquanto a maior Média foi obtida na região Sudeste ( 28,5 ).
Quanto à variabilidade das notas, o Desvio padrão de todo o Brasil foi 32,1.
O maior Desvio padrão foi obtido na região Sudeste ( 33,8 ), enquanto o menor foi obtido na região Norte ( 27,0 ).
A nota Máxima, 100,0 pontos, foi alcançada por, pelo menos, um aluno de todas as regiões.
A Mediana do Brasil, como um todo, foi cinco, a mesma para a região Sul.
As regiões Norte, Nordeste e Centro-Oeste tiveram Mediana zero, significando que, pelo menos, a metade dos alunos em cada uma dessas regiões tirou nota zero nessa questão.
A região Sudeste obteve a maior Mediana, 15,0.
A nota Mínima ( 0,0 ) também foi a mesma em todas as regiões do Brasil. 6.3.2.2 - Comentários sobre as respostas à Questão Discursiva 3 Sob a perspectiva das Diretrizes Curriculares Nacionais para o Curso de Graduação em Ciências Contábeis, a questão abordou conteúdos que, em geral, são trabalhados em No Gráfico 6.11, mostra -se a distribuição das notas na questão discursiva 3, do Componente de Conhecimento Específico da área de Ciências Contábeis.
Sem se levar em consideração que acima de 32 \% dos estudantes deixaram a questão em branco e que quase 16 \% receberam nota zero, a moda dessa distribuição seria o intervalo ( 10 ; 20 ], com alguns máximos locais. 6.3.2.2 - Comentários sobre as respostas à Questão Discursiva 3 Sob a perspectiva das Diretrizes Curriculares Nacionais para o Curso de Graduação em Ciências Contábeis, a questão abordou conteúdos que, em geral, são trabalhados em disciplinas de Contabilidade de Custos e Contabilidade Gerencial.
Mais especificamente foram abordados os conceitos de margem de contribuição unitária, margem de contribuição total, despesas e custos fixos e variáveis, método de custeio variável e ponto de equilíbrio.
Todos esses conceitos são relativos à análise da relação custo-volume-lucro.
As informações necessárias para responder aos itens do problema eram claras e suficientes.
Além disso, o enunciado apresentava diversos conceitos, relacionados ao custeio variável, que auxiliavam a resolução da questão, especificamente o conceito de margem de contribuição, cobrado no item ‘ a ’.
A questão não requeria conteúdos aprofundados ou desenvolvimento de análises apuradas.
Era bastante simples, com um nível de profundidade básico, pela ênfase no cálculo de valores, a partir da aplicação de uma lógica conceitual sobre análise de custo-volume-lucro.
Assim, as respostas envolviam apenas o cálculo dos valores correspondentes à margem de contribuição, ao lucro pelo custeio variável e ao ponto de equilíbrio, sem solicitação de considerações sobre os valores apurados.
As respostas encontradas foram, em geral, compatíveis com o padrão, exceto quanto à exigência de detalhamento dos cálculos realizados que, em muitos casos, não foram apresentados.
No item ' a, ' muitas respostas apresentaram apenas o valor da margem de contribuição unitária dos produtos, o que indicou a não compreensão da necessidade de demonstrar o cálculo.
Outras respostas trouxeram a margem de contribuição unitária como percentual da receita de vendas.
Outro equívoco bastante comum encontrado nas respostas foi erro de cálculo.
Mesmo indicando a \textbf{fórmula} correta de apuração da margem de contribuição unitária, muitos erraram as operações, que envolviam soma e subtração.
Foram verificadas muitas respostas incorretas por não considerarem que, na apuração da margem de contribuição unitária do produto, deveriam ser subtraídos do preço de venda os custos e despesas variáveis e não apenas um dos dois.
O item ' b ' solicitou a apuração do lucro da empresa pelo custeio variável, considerando um dado volume de vendas para dois produtos : X, com 30 unidades e Z, com 40 unidades.
Após a apuração da margem de contribuição total, deveria ser excluído o valor dos custos fixos, porém em muitas respostas o custo fixo foi excluído da margem de contribuição de cada produto individualmente.
Houve casos, também, em que o valor do custo fixo foi rateado entre os dois produtos, a critério do respondente.
Esses erros demonstram desconhecimento dos estudantes quanto à forma de apuração do resultado pelo método do custeio variável.
Também foram observados muitos casos de respostas apenas com o valor final do lucro da empresa ou de cada um dos produtos separadamente.
Os erros nas operações de soma e disciplinas de Contabilidade de Custos e Contabilidade Gerencial.
Mais especificamente foram abordados os conceitos de margem de contribuição unitária, margem de contribuição total, despesas e custos fixos e variáveis, método de custeio variável e ponto de equilíbrio.
Todos esses conceitos são relativos à análise da relação custo-volume-lucro.
As informações necessárias para responder aos itens do problema eram claras e suficientes.
Além disso, o enunciado apresentava diversos conceitos, relacionados ao custeio variável, que auxiliavam a resolução da questão, especificamente o conceito de margem de contribuição, cobrado no item ‘ a ’.
A questão não requeria conteúdos aprofundados ou desenvolvimento de análises apuradas.
Era bastante simples, com um nível de profundidade básico, pela ênfase no cálculo de valores, a partir da aplicação de uma lógica conceitual sobre análise de custo-volume-lucro.
Assim, as respostas envolviam apenas o cálculo dos valores correspondentes à margem de contribuição, ao lucro pelo custeio variável e ao ponto de equilíbrio, sem solicitação de considerações sobre os valores apurados.
As respostas encontradas foram, em geral, compatíveis com o padrão, exceto quanto à exigência de detalhamento dos cálculos realizados que, em muitos casos, não foram apresentados.
No item ' a, ' muitas respostas apresentaram apenas o valor da margem de contribuição unitária dos produtos, o que indicou a não compreensão da necessidade de demonstrar o cálculo.
Outras respostas trouxeram a margem de contribuição unitária como percentual da receita de vendas.
Outro equívoco bastante comum encontrado nas respostas foi erro de cálculo.
Mesmo indicando a \textbf{fórmula} correta de apuração da margem de contribuição unitária, muitos erraram as operações, que envolviam soma e subtração.
Foram verificadas muitas respostas incorretas por não considerarem que, na apuração da margem de contribuição unitária do produto, deveriam ser subtraídos do preço de venda os custos e despesas variáveis e não apenas um dos dois.
O item ' b ' solicitou a apuração do lucro da empresa pelo custeio variável, considerando um dado volume de vendas para dois produtos : X, com 30 unidades e Z, com 40 unidades.
Após a apuração da margem de contribuição total, deveria ser excluído o valor dos custos fixos, porém em muitas respostas o custo fixo foi excluído da margem de contribuição de cada produto individualmente.
Houve casos, também, em que o valor do custo fixo foi rateado entre os dois produtos, a critério do respondente.
Esses erros demonstram desconhecimento dos estudantes quanto à forma de apuração do resultado pelo método do custeio variável.
Também foram observados muitos casos de respostas apenas com o valor final do lucro da empresa ou de cada um dos produtos separadamente.
Os erros nas operações de soma e subtração foram frequentes, o que resultou na apuração errada do lucro pelo custeio variável, mesmo conhecendo a lógica de cálculo.
O item ' c ' solicitou a apuração do ponto de equilíbrio contábil, obtido a partir da divisão dos custos fixos pela margem de contribuição unitária, no caso do produto X.
Foram observadas muitas respostas que apresentaram a \textbf{fórmula} de cálculo do ponto de equilíbrio contábil sem chegar ao valor solicitado.
Também foi comum o respondente apenas indicar o valor como resposta, sem apresentar o cálculo.
Neste item, também foram observadas muitas respostas erradas em decorrência de erro de cálculo, que requeria uma divisão simples.
Merece destaque o fato de a resposta ao item ’ c ’ depender da apuração correta da margem de contribuição unitária do produto, solicitada no item ' a ', tornando -os interdependentes.
A linguagem adotada pelos alunos foi adequada do ponto de vista dos termos específicos da área, inclusive porque o enunciado da questão \textbf{ajudou} na definição dos conceitos relacionados ao assunto abordado.
O desempenho geral na Questão 3 ficou abaixo do esperado, com média 24,4 e um percentual de notas zero superior a 18 \% do total corrigido, enquanto as respostas em branco, anuladas ou desconsideradas ultrapassaram 30 \% do total de alunos que fizeram o Exame.
Essa questão, que não requer conteúdo atualizado, mas apenas conhecimentos que são ministrados em qualquer curso de Ciências Contábeis, teve em torno de 40 \% das respostas fracas ( notas de zero a 30 ), percentual relativo aos que apresentaram alguma resposta para a questão e não acertaram nem, ao menos, um dos três itens solicitados.
Pouco menos de 20 \% das respostas corrigidas foram consideradas medianas ( notas de 35 a 70 ), em que os alunos conseguiram acertar de forma completa, pelo menos, um dos três itens da questão.
As respostas consideradas boas ( notas de 75 a 100 ) foram alcançadas por cerca de 10 \% dos alunos e apresentaram, de forma satisfatória, a maior parte do que se exigia.
As respostas encontradas revelam falta de domínio, por boa parte dos alunos, acerca de conteúdo básico na formação do bacharel em Ciências Contábeis.
Além de falta de domínio conceitual, muitos alunos demonstraram ter dificuldade na realização de operações matemáticas básicas.
Tal dificuldade levou, inclusive, a várias manifestações de insatisfação pela impossibilidade de uso da calculadora no exame, o que era completamente desnecessário pela facilidade dos cálculos exigidos.
Apesar desse relato, entre as três questões discursivas de Conhecimentos Específicos de Ciências Contábeis do Enade-2018, essa foi a questão que teve o menor percentual de respostas em branco, apesar de alto ( em torno de 30 \% ), e a maior nota média. 6.3.2.3 - Análise da Questão Discursiva 4 do Componente de Conhecimento Específico Na Tabela 6.19, constam as informações relativas à questão 4 do conjunto de questões do Componente de Conhecimento Específico.
O desempenho dos estudantes nessa questão foi inferior ao das questões 3 e 5.
A Média geral do Brasil foi 8,1, sendo a maior Média registrada na região Sudeste ( 9,2 ), e a menor, na região Norte ( 5,2 ).
A nota Máxima ( 100,0 ) foi atingida por, pelo menos, um concluinte de todas as regiões.
A Mediana em todo o Brasil foi zero, o mesmo valor foi obtido nas cinco Grandes Regiões, indicando que, pelo menos, 50 \% dos participantes tiraram nota zero.
Também foi zero a nota Mínima de todas as regiões.
No Gráfico 6.12, está representada a distribuição de notas da questão discursiva 4, do Componente de Conhecimento Específico.
Essa distribuição tem moda na classe de estudantes que deixaram a questão em branco, seguida por aqueles que, tendo respondido à questão, obtiveram nota zero.
A soma das frequências destes dois intervalos supera 80 \% dos respondentes.
O intervalo ( 20 ; 30 ] caracteriza -se como um máximo local, com cerca de 9 \% do total. 6.3.2.4 Comentários sobre as respostas à Questão Discursiva 4 A Questão 4 abordou conteúdos de Contabilidade Societária que, em geral, são trabalhados em disciplinas de Contabilidade Intermediária e Contabilidade Avançada e estão presentes nas Diretrizes Curriculares Nacionais para o Curso de Graduação em Ciências Contábeis.
Avaliava competência relativa à interpretação e aplicação de normas contábeis para identificar a necessidade de revisão no valor de um item patrimonial.
O conteúdo avaliado pode ser considerado de nível intermediário de dificuldade, mas a questão foi apresentada de forma simples e com um enunciado que trazia elementos para auxiliar a resolução dos dois itens requeridos pelo comando.
Vale destacar, inclusive, que, a partir do enunciado, era possível responder a parte do item ‘ a ’, mesmo sem conhecer exatamente os conceitos relativos ao assunto.
A questão pode ser considerada fácil, por não exigir análises complexas, mas apenas aplicação de conceitos e domínio do método de lançamentos contábeis.
Até mesmo o valor da depreciação acumulada do ativo foi informado no enunciado, sendo desnecessário esse cálculo auxiliar.
Apesar de o enunciado estar claro e não dar margem para interpretações diferentes, muitas respostas indicaram entendimento equivocado, como se fosse uma questão de depreciação e não de redução ao valor recuperável de ativos.
Em algumas respostas, os alunos cometeram erros conceituais preocupantes.
Comparativamente às demais questões discursivas de Conhecimentos Específicos, essa foi a questão de conteúdo mais atualizado e aprofundado, uma vez que abordava uma norma contábil adotada há cerca de 10 anos no Brasil, mais especificamente para as empresas de capital aberto e sociedades anônimas.
Essas empresas são a minoria entre os milhares de empresas em atividade no Brasil e não fazem parte da realidade do universo profissional da maioria dos contadores brasileiros.
O item ' a ' solicitava o valor a ser apresentado como saldo do ativo imobilizado ao final do exercício, com a devida fundamentação.
O valor do saldo do ativo imobilizado ao final do exercício é de R\$40.000,00, que foi indicado corretamente em muitos casos, porém sem a devida fundamentação.
Como já comentado, era possível acertar parcialmente o item ' a ' ( valor do saldo do ativo imobilizado ao final do exercício ), sem nenhum conhecimento do conceito de redução ao valor recuperável, porque ( i ) o valor já estava explícito na questão, e ( ii ) esse mesmo valor pode ser obtido pela subtração entre o custo de aquisição e o valor residual do ativo.
Foram frequentes respostas que se restringiram ao item ‘ a ’.
O item ' b ' solicitava que fossem explicitados os lançamentos contábeis a serem realizados para evidenciar o valor correto do ativo imobilizado.
Foram observados muitos casos de lançamentos com contas equivocadas ou com valor equivocado, além de ausência de justificativa para o lançamento.
Em relação a esse item, observou -se que os alunos têm dificuldade na identificação de origem e aplicação de recursos, que caracteriza o método das partidas dobradas.
Foram poucos os casos de lançamentos apresentados adequadamente.
Observou -se, também, uma alta frequência de respostas incompletas.
Quanto ao item ' a ', foi mais frequente o estudante apresentar apenas o valor do saldo do ativo imobilizado, sem nenhuma fundamentação acerca do valor indicado ; e no que tange ao item ' b ', as repostas corretas mais frequentes foram de estudantes que apresentaram apenas o lançamento contábil em diário ou em razonetes.
Nas tentativas de resposta à questão, a linguagem adotada pelos alunos foi adequada do ponto de vista dos termos específicos da área, inclusive porque o enunciado da questão \textbf{ajudou} na definição dos conceitos relacionados ao assunto abordado.
Os erros mais encontrados parecem indicar que os alunos confundem redução ao valor recuperável de ativos com depreciação, que é um conteúdo básico em Contabilidade.
A A grande maioria das respostas diferentes do esperado eram incompatíveis com o padrão de resposta e sem possibilidade de ser consideradas corretas.
O desempenho geral das respostas da Questão 4 também foi muito fraco.
Mais de 50 \% dos estudantes que tentaram responder à questão obtiveram notas zero, e a média foi 8,1, menor que que 10 numa escala de zero a 100.
As respostas em branco superaram 32 \%, e houve uma alta quantidade de respostas anuladas ou desconsideradas entre as corrigidas.
Esses dados, associados a uma média muito baixa, representam um resultado preocupante.
Em torno de 92 \% das respostas corrigidas foram consideradas fracas ( notas de zero a 30 ), nas quais os alunos não conseguiram responder, de forma satisfatória e integral, a nenhum dos dois itens.
Cerca de 6 \% das respostas corrigidas foram consideradas medianas ( notas de 35 a 70 ), notas atribuídas a alunos que conseguiram apresentar o conteúdo requerido de forma satisfatória em pelo menos um dos itens da questão.
As respostas com notas consideradas boas ( notas de 75 a 100 ) foram atribuídas a apenas 2 \% dos alunos que responderam à questão, por apresentarem de forma satisfatória boa parte do que era solicitado nos dois itens da questão ou acertá -los integralmente.
As respostas encontradas revelam um significativo desconhecimento dos alunos acerca do tema abordado, tópico de nível intermediário na formação do bacharel em Ciências Contábeis.
A falta de domínio teórico e conceitual, associada a uma leitura apressada do comando, talvez explique o fato de grande parte das respostas não apresentarem a fundamentação para que o valor do ativo imobilizado ao final do exercício fosse reduzido.
O resultado pode estar associado à desatualização do ensino em relação às normas contábeis aplicáveis às empresas de capital aberto e sociedades anônimas em geral no Brasil. 6.3.2.5 Análise da Questão Discursiva 5 do Componente de Conhecimento Específico Na Tabela 6.20, constam as informações relativas à questão 5 do conjunto do Componente de Conhecimento Específico.
O desempenho dos estudantes de todo o Brasil nesta questão foi inferior ao da questão 3 e superior ao da questão 4.
A nota Média dos estudantes de todo o Brasil foi 11,3.
A maior Média foi registrada na região Sudeste ( 12,8 ), enquanto a menor Média foi registrada na região Norte ( 9,5 ).
Quanto à variabilidade das notas, o Desvio padrão dos alunos do Brasil, como um todo, foi 20,5.
Enquanto o maior desvio foi encontrado na região Sudeste ( 21,9 ), o menor foi encontrado na região Norte ( 18,7 ).
A Mediana para o Brasil e em todas as regiões foi zero, associada ao excessivo número de notas zero, como mostrado no Gráfico 6.13, indicando que mais de metade dos alunos receberam esta nota, em cada região.
Para o conjunto de alunos de Ciências Contábeis do Brasil, a nota Máxima foi 100,0, obtida, também, em todas as regiões.
A nota Mínima do Brasil foi zero, o mesmo valor encontrado em todas as regiões.
No Gráfico 6.13, é representada a distribuição das notas da questão discursiva 5 do Componente de Conhecimento Específico.
Essa distribuição, como todas as relativas às questões discursivas dessa parte da prova, também tem como modal o intervalo correspondente ao dos que deixaram a questão em branco, com 41,5 \% do total.
Somando -se a frequência desse intervalo com a daqueles que obtiveram nota zero chega -se próximo a 70 \% dos participantes. 6.3.2.6 Comentários sobre as respostas à Questão Discursiva 5 A Questão 5 abordava conteúdos de Auditoria, em geral trabalhados a partir da segunda metade do curso, momento em que os estudantes já deveriam ter desenvolvido o domínio da lógica contábil e poderiam associar as informações contábeis no contexto do trabalho do auditor, bem como suas implicações para a continuidade das operações de uma entidade.
Tais conhecimentos estão em consonância com as Diretrizes Curriculares Nacionais para o Curso de Graduação em Ciências Contábeis.
No enunciado da questão, contextualizava -se a situação da continuidade operacional de uma empresa sob a perspectiva de crises econômicas, principalmente nos casos em que a empresa depende do financiamento do capital de terceiros.
Tanto na elaboração das demonstrações contábeis, quanto na atividade de auditoria das mesmas, o pressuposto da continuidade é importante, pois afeta a mensuração dos elementos patrimoniais ( ativos e passivos ) e a análise que se faz de tais informações.
A questão exigia do aluno o julgamento de situações que podem ser encontradas nas atividades de um profissional de auditoria, que tem a responsabilidade de emitir um parecer sobre a adequação das informações contábeis para representar a realidade operacional de uma empresa.
O comando estava subdividido em três itens.
O item ' a ' solicitava que o aluno citasse os tipos de opinião que o auditor independente pode emitir em seu relatório de auditoria.
O item ' b ' abordava a influência da dúvida substancial quanto à continuidade operacional na formação da opinião do auditor.
Por fim, o item ' c ' tratava do ceticismo profissional e sua influência na decisão do auditor em emitir uma opinião acerca da continuidade operacional de uma empresa.
Enquanto o primeiro item requeria a simples indicação de tipos de pareceres que podem ser emitidos, os demais itens eram mais subjetivos e requeriam competências de análise e reflexão por parte do respondente.
A questão foi apresentada em linguagem clara e continha informações suficientes para o estudante responder aos três itens, embora tivesse poucos elementos no enunciado que auxiliassem diretamente na formulação das respostas.
A Questão 5, em seu conjunto, pode ser considerada de nível de dificuldade médio, porém mais difícil do que as demais questões discursivas de Conhecimentos Específicos.
Um dos itens exigia apenas a memorização de nomenclatura dos tipos de opinião, enquanto os outros dois requeriam a análise de um fato hipotético apresentado no enunciado e o uso de conhecimento sobre normas e procedimentos de auditoria para tomada de decisão em uma atividade de auditoria em uma empresa.
O mais frequente foi o estudante responder apenas ao item ' a ', que pedia para indicar os tipos de opinião que o auditor independente pode emitir em seu relatório.
Não foram poucos os casos em que a linguagem adotada não era adequada sob a perspectiva dos termos específicos da área, constantes nas normas e procedimentos de auditoria emitidos pelo CFC.
Os alunos indicavam os tipos de opinião utilizando outros termos não condizentes com o previsto nas normas de auditoria, tais como opinião favorável e desfavorável ou contrária ; opinião regular, com ou sem ressalva, e irregular ; opinião positiva e negativa ; opinião ou parecer limpo ; e ainda impossibilidade ou incapacidade de emitir opinião.
Embora alguns desses termos possam ser usados no jargão dos auditores, não condizem com a nomenclatura técnica constante das normas de auditoria.
Com relação aos itens ' b ' e ' c ', além de menos tentativas de responder, muito poucos conseguiram se aproximar do padrão de respostas, o que resultou em um desempenho muito fraco.
As poucas respostas corrigidas não eram aceitáveis, demonstrando falta de conhecimento do assunto abordado ou, ainda, incapacidade de articulação entre os conceitos e procedimentos de auditoria e a situação hipotética apresentada no enunciado da questão.
Ressalta -se que não se trata de conteúdo novo ou avançado do curso de formação em Ciências Contábeis.
Foram observados muitos casos em que os alunos demonstravam não compreender o conceito de continuidade operacional, nem o papel e a responsabilidade do auditor na emissão do relatório com a sua opinião.
O percentual de notas zero das respostas da Questão 5 foi muito alto, em torno de 35 \% do total corrigido, enquanto as respostas em branco, anuladas ou desconsideradas alcançaram mais de 45 \% do total de alunos que fizeram o Exame.
Essa questão, de conteúdo intermediário na formação de um bacharel em Ciências Contábeis, teve uma média de pontuação muito baixa, 11.3, numa escala de zero a 100, o que representa um resultado preocupante.
Em torno de 90 \% das respostas foram consideradas fracas ( notas de zero a 30 ), nas quais os alunos não conseguiram responder de forma correta e integral sequer ao item ‘ a ’.
Em torno de 8 \% das respostas foram consideradas medianas ( notas de 35 a 70 ), respostas satisfatórias em pelo menos um dos itens da questão, frequentemente relativo ao primeiro deles.
E as respostas consideradas boas ( notas de 75 a 100 ) foram apenas 2 \% do total de alunos que fizeram a prova.
Em geral, o desempenho dos estudantes nessa questão foi considerado muito fraco, o que sugere uma lacuna entre o que foi exigido no Exame e os conhecimentos desenvolvidos pelos alunos, ao longo de seus cursos de formação em relação à matéria de Auditoria.
A partir da correção da Questão 05, considera -se que as respostas sinalizam uma formação deficiente do ponto de vista de domínio teórico e conceitual do assunto, além de baixa capacidade de articulação.
Algumas respostas indicavam uma formação voltada para aspectos tradicionais da área contábil, como questões legais e tributárias.
A capacidade de redação e expressão de ideias também se relevou deficiente.
As respostas, em geral, eram mal elaboradas, sem argumentos adequados, com textos simplórios e incompatíveis com o nível de formação esperado.
Entre as questões de conhecimentos específicos de Ciências Contábeis do Enade-2018, esta foi a que teve o maior percentual de respostas em branco ( 42 \% ) e, mais uma vez, uma média muito baixa, decorrente do provável desconhecimento sobre o tema e também da reduzida capacidade de redação e expressão dos alunos. 6.3.3 Considerações Finais As questões de conhecimentos específicos abordadas no Enade-2018 versaram sobre os temas : Contabilidade de Custos e Gerencial ( Método de Custeio Variável e Relação Custo-volume-lucro ), Redução ao Valor Recuperável de Ativos e Auditoria.
Em geral, o nível de dificuldade das questões foi considerado de fácil a médio, com pouca exigência de raciocínio analítico ou articulação com outros temas.
Porém, o desempenho observado revelou indícios de deficiência na formação dos concluintes, até mesmo por se tratar de conteúdos básicos do curso, que não exigem atualização normativa e teórica.
Além disso, as respostas redigidas pelos estudantes indicam baixa capacidade de expressão de ideias, mesmo quando elementos auxiliares estavam no próprio enunciado.
Ao analisar as três questões de conhecimentos específicos, verifica -se que as notas das questões ficaram muito abaixo de 50 \% da pontuação máxima, com uma quantidade expressiva de respostas em branco e respostas com nota zero.
Esses resultados tanto podem sugerir a existência de deficiências na formação dos alunos, como também podem estar indicando a falta de comprometimento com a realização do exame.

% matched lemmas: ajudar, fórmula, laranja
\end{textsample}
